\documentclass[11pt,a4paper]{article}
\usepackage[margin=2.9cm]{geometry}
\usepackage{amsmath,amssymb,amsthm,mathrsfs}
\usepackage[T1]{fontenc}
\usepackage{microtype}
\usepackage{booktabs}
\usepackage[colorlinks=true,linkcolor=blue,citecolor=blue]{hyperref}

\theoremstyle{plain}
\newtheorem{theorem}{Theorem}[section]
\newtheorem{proposition}[theorem]{Proposition}
\newtheorem{lemma}[theorem]{Lemma}
\newtheorem{corollary}[theorem]{Corollary}
\newtheorem{hypothesis}[theorem]{Hypothesis}
\theoremstyle{definition}
\newtheorem{definition}[theorem]{Definition}
\newtheorem{remark}[theorem]{Remark}
\newtheorem{srcfact}[theorem]{Source Fact}

\newcommand{\R}{\mathbb{R}}
\newcommand{\Z}{\mathbb{Z}}
\newcommand{\E}{\mathbb{E}}
\newcommand{\A}{\mathfrak{A}}
\newcommand{\G}{\mathcal{G}}
\newcommand{\Q}{\mathcal{Q}}
\newcommand{\HH}{\mathcal{H}}
\newcommand{\Om}{\Omega}
\newcommand{\ip}[2]{\langle #1,#2\rangle}
\DeclareMathOperator{\supp}{supp}
\DeclareMathOperator{\dist}{dist}
\DeclareMathOperator{\spec}{spec}
\DeclareMathOperator{\Span}{span}

\title{Removal of the spatial cutoff for weakly coupled $P(\varphi)_2$:\\
norm convergence of the ground-state net on the quasi-local algebra\\[4pt]
\large (M6, v2 --- audited repair)}
\author{}
\date{}

\begin{document}
\maketitle

\begin{abstract}
Let $H(g)$ be the spatially cutoff $P(\varphi)_2$ Hamiltonian, $\Om_g$ its ground state, and
$\omega_g$ the induced state on the time-zero quasi-local algebra $\A$. Glimm and Jaffe (1970)
constructed an infinite-volume vacuum as a $w^*$ \emph{limit point} of space-averaged $\omega_g$;
the cited construction does not prove convergence of the net itself. Guerra--Rosen--Simon
(1975) formulated analogous Euclidean local-convergence questions as their Problems~1--2.
We prove: in the Glimm--Jaffe--Spencer weak-coupling region
$0\le\lambda/m_0^2<\varepsilon_{\mathscr P}$, the \emph{full net} $(\omega_g)$ --- no subnets, no
averaging --- converges in norm on every local algebra,
\[
  \lim_{g\in\G}\bigl\|(\omega_g-\omega_\infty)|_{\A(O)}\bigr\|=0 ,
\]
with an explicit super-polynomial rate in $\dist(O,\{g<1\})$; the limit is locally normal,
translation invariant, $\Z_2$-invariant for even $\mathscr P$, and is a ground state of the
Glimm--Jaffe dynamics $(\A,\alpha)$ in the spectral sense. The proof is Hamiltonian: a
local-perturbations-perturb-locally estimate built on an exact spectral filter and the exact
light cone makes the net Cauchy; its two hypotheses --- a cutoff-uniform gap
$\spec H(g)\subseteq\{0\}\cup[m_1,\infty)$ and a uniform local interaction moment --- are derived
from the two printed cluster-expansion theorems of Glimm--Jaffe--Spencer (Ann.\ of Math.\ 100
(1974), Thms.~1.1.7 and 1.1.8), whose cutoff-uniformity is both printed (p.~629) and audited.
The proof also imports, and explicitly identifies below, the standard cutoff-construction
results on self-adjointness and the simple positive ground state, Feynman--Kac--Nelson transfer,
finite-volume stability and hypercontractivity, domain smoothing and a common core, finite
propagation and patching, and covariance. No Euclidean uniqueness theorem, global Markov
property for the infinite-volume measure, or reverse Osterwalder--Schrader reconstruction is
used. We do not claim uniqueness of $\omega_\infty$ among all locally normal ground states, a
sharp exponential rate, or anything outside the weak-coupling region.
\end{abstract}

\tableofcontents

\section{Introduction}\label{sec:intro}

\subsection{The problem}

The $\lambda\varphi^4_2$ (more generally $P(\varphi)_2$) model was constructed by Glimm and
Jaffe in the series I--IV \cite{GJ1,GJ2,GJ3,GJ4}. Paper III builds the physical vacuum as
follows: the finite-volume ground states $\Om_g$ of the cutoff Hamiltonians $H(g)$ induce states
$\omega_g$ on the quasi-local algebra $\A$; after a space-translation averaging, a $w^*$-compactness
argument produces a \emph{limit point} $\omega$, which is shown to be locally Fock
\cite[Thms.~2.1--2.3]{GJ3}. Whether the net $(\omega_g)$ \emph{converges} --- so that the vacuum
is canonical rather than selected by a limit procedure --- was not proved there; the Expositions volume states plainly that
the $g\to1$ limit of the vacuum expectation values ``has not been proved''
\cite[\S4]{GJExp}. Guerra, Rosen and Simon list analogous \emph{Euclidean} local $L^p$ and
finite-volume Gibbs-state convergence statements as Problems~1--2 of \cite[pp.~125--126]{GRS75};
their note added on p.~128 records Newman's solution of the small-coupling $p=1$ versions.
(Some later surveys state the Hamiltonian cutoff limit as a convergence; the cited primary
construction itself is unambiguously a limit-point construction.)

On the Euclidean side, convergence of \emph{Schwinger functions} at weak coupling follows from
the cluster expansion of Glimm--Jaffe--Spencer \cite{GJS}. Albeverio--H\o egh-Krohn--Zegarlinski
prove Euclidean Gibbs uniqueness under their stated hypotheses \cite{AHZ}; Bauerschmidt--Dagallier--
Weber prove uniqueness of the invariant measure for the $\Phi^4_2$ stochastic-quantization
equation above the critical temperature \cite{BDW}. Neither statement is, by itself, a theorem
about norm convergence of this Hamiltonian ground-state cutoff net. The cited primary
construction does not prove that local-norm convergence, and no such theorem is identified in
the references reviewed here. The Hamiltonian-renormalisation programme of Rodriguez
Zarate--Thiemann cited here is explicitly restricted to finite volume \cite{RZT}.

\subsection{The result}

\begin{quote}
\emph{In the GJS weak-coupling region, the full net $(\omega_g)_{g\in\G}$ converges in norm on
every local algebra; the limit is locally normal, translation and (for even $\mathscr P$)
$\Z_2$ invariant, and is a ground state of the Glimm--Jaffe dynamics. An explicit rate holds:
$\|(\omega_g-\omega_\infty)|_{\A(O)}\|\le\Psi_{m_1}(d)$ for $g\equiv1$ on $O_d$, with
$\Psi_{m_1}$ decaying faster than every inverse power.}
\end{quote}

See Theorem~\ref{thm:MAIN}. The logical structure is deliberately Hamiltonian and
compactness-free:

\begin{center}
\begin{tabular}{ccc}
\toprule
step & content & where proved\\
\midrule
(i) & printed GJS cluster/moment theorems, cutoff-uniform & \cite{GJS} + audit \S\ref{sec:sources}\\
(ii) & uniform local moment (M$_2$) & \S\ref{sec:M2}\\
(iii) & uniform gap (G): $\spec H(g)\subseteq\{0\}\cup[m_1,\infty)$ & \S\ref{sec:G}\\
(iv) & LPPL/Cauchy: (G)+(M$_2$) $\Rightarrow$ full-net convergence + rate & \S\ref{sec:lppl}\\
(v) & the limit is an $\alpha$-ground state; invariances & \S\ref{sec:ground}\\
\bottomrule
\end{tabular}
\end{center}

Steps (ii)--(iii) convert Euclidean input to Hamiltonian statements through slab
Feynman--Kac--Nelson transfers in which every limit is taken at fixed $g$; no uniform gap is
presupposed anywhere in their proofs (the gap of step (iii) is an output). Step (iv) replaces
the compactness of \cite{GJ3} by a Cauchy argument: the same gap that yields clustering also
yields a local-perturbations-perturb-locally estimate, which bounds
$\|(\omega_g-\omega_{g'})|_{\A(O)}\|$ whenever both cutoffs equal $1$ on $O_d$. Step (v) is the
exact spectral identity of the companion note M0: the ground-state property of the limit costs
an identity, not an estimate.

\subsection{What is not claimed}

(a) \emph{Uniqueness}: Theorem~\ref{thm:MAIN} identifies the limit of \emph{this} cutoff family;
it does not show that every locally normal ground state of $(\A,\alpha)$ equals $\omega_\infty$.
(b) \emph{Sharp rate}: in the exactly solvable quadratic model the true rate is
$e^{-2m_1d}$ (companion note M2); $\Psi_{m_1}$ is super-polynomial but not exponential, and this
is intrinsic to the filter method used. (c) \emph{Beyond weak coupling}: under the explicit
phase-classification hypotheses stated in Corollary~\ref{cor:nobreak}, a cutoff-uniform gap is
incompatible with that two-phase structure. Thus an extension to such a broken phase would have
to abandon the uniform-gap input; no claim is made outside
$0\le\lambda/m_0^2<\varepsilon_{\mathscr P}$. (d) \emph{Wavelet uniformity}: uniformity in a
lattice/wavelet resolution is a separate theorem (Part II) and is not addressed here.

\subsection*{Provenance}

This is the v2 audited repair: all components are stated and proved in this document. The LPPL machinery
of \S\ref{sec:lppl} originates in the companion files \texttt{m4/lppl.tex} and
\texttt{m4/remaining.tex} and has been re-verified line by line during the merge (two sign/constant
slips in the v0 pointers were corrected: the filter transform is $+E^{-1}$, the one-path bound
carries no factor $\tfrac12$). This version additionally repairs the source pinpoints, the
dimensionless scaling reduction, signed and unbounded FKN passages, the unit-cell count, the
negative-time norm ratio, the exponential-moment prefactor, the spectral-projection range, and
the hypotheses of the symmetry-breaking corollary. External inputs are the printed theorems of
\cite{GJS} (with the uniformity audit summarised in Remark~\ref{rem:audit}) and the explicitly
cited construction results from \cite{SimonBook}, \cite{GRS75}, and \cite{GJExp}; companion notes
are used only for provenance and extended discussion. Gate R2 (the norm prefactor of
\cite[(1.1.7)--(1.1.8)]{GJS}) is closed by direct inspection of the page images: $K_1$ occurs once
and the local weight is $K_2^{2N(\Delta)}$.

\section{Setting and standing assumptions}\label{sec:setting}

\subsection{The model}

$\HH=L^2(\mathcal S'(\R),d\phi_0)$ is $Q$-space over the time-zero free field of mass $m_0>0$:
$d\phi_0$ is Gaussian with covariance $C=(2\mu)^{-1}$, $\mu=(-\partial_x^2+m_0^2)^{1/2}$. Weyl
operators $W(f,h)=\exp i(\varphi_0(f)+\pi_0(h))$ generate the local net
\[
  \A(O)=\{W(f,h):f,h\in C_c^\infty(O)\ \text{real}\}'',\qquad
  \A=\overline{\textstyle\bigcup_O\A(O)}^{\;\|\cdot\|} .
\]
Let $\mathscr P(\xi)=\sum_{n=0}^{2p'}a_n\xi^n$, $a_{2p'}>0$, and for real $h\in C_c^\infty$
\[
  V(h):=\lambda\int h(x)\,{:}\mathscr P(\varphi_0(x)){:}_C\,dx ,
\]
a self-adjoint multiplication operator in $Q$-space, affiliated with $\A(U)$ for
$U\supseteq\supp h$, finitely additive on common domains ($V(\sum h_j)\xi=\sum V(h_j)\xi$).

\[
  \G:=\{g\in C_c^\infty(\R):0\le g\le1,\ g\equiv1\text{ on a nonempty bounded interval}\},
\]
with $g\preccurlyeq g'$ exactly when
$\{g=1\}^\circ\subseteq\{g'=1\}^\circ$. For a bounded interval $O$ and $r\ge0$, set
\[
 O_0:=O,\qquad O_r:=\{x\in\R:\dist(x,O)<r\}\quad(r>0),\qquad
 d(g,O):=\dist\bigl(O,\R\setminus\{g=1\}^\circ\bigr).
\]
\[
  K(g):=H_0+V(g),\quad E(g):=\inf\spec K(g),\quad H(g):=K(g)-E(g),
\]
with ground state $\Om_g$, $\omega_g:=\ip{\Om_g}{\cdot\,\Om_g}$,
$\beta^g_t=\mathrm{Ad}\,e^{itH(g)}$.

\subsection{Standing assumption (weak coupling)}

Fix source-admissible constants $M_*>0$ and $\lambda_*>0$ such that the
Glimm--Jaffe--Spencer expansion is valid at bare mass $M_*$ whenever
$0\le\lambda'<\lambda_*$. Define
\[
  \varepsilon_{\mathscr P}:=\frac{\lambda_*}{M_*^2}.
\]
Throughout, $0\le\lambda/m_0^2<\varepsilon_{\mathscr P}$. After transport back to the original
coordinates, $m_1>0$ and $C_7,C_8,K_1,K_2$ denote the constants in GJS
Theorems 1.1.7--1.1.8 and their observable norms; only their positivity or finiteness and their
cutoff independence will be used.

\begin{proposition}[dimensionless weak-coupling reduction]\label{prop:scaling}
The preceding inequality implies the source smallness hypothesis. More precisely, the GJS
unitary dilation with $s=m_0/M_*$ sends
\[
  (m_0,\lambda)\longmapsto
  (s^{-1}m_0,s^{-2}\lambda)
  =\left(M_*,\frac{M_*^2}{m_0^2}\lambda\right),
\]
and the transformed coupling is strictly smaller than $\lambda_*$.
\end{proposition}

\begin{proof}
For $m>0$, write the Euclidean covariance, as an identity of tempered distributions, as
\[
  C_m(z)=\frac1{(2\pi)^2}\int_{\R^2}\frac{e^{ip\cdot z}}{p^2+m^2}\,dp .
\]
Since $sM_*=m_0$, the substitution $q=sp$ gives, without an omitted power of $s$,
\begin{align*}
  C_{M_*}(sz)
  &=\frac1{(2\pi)^2}\int_{\R^2}
      \frac{e^{ip\cdot sz}}{p^2+M_*^2}\,dp \\
  &=\frac1{(2\pi)^2}\int_{\R^2}
      \frac{s^{-2}e^{iq\cdot z}}{s^{-2}q^2+M_*^2}\,dq
    =C_{m_0}(z).
\end{align*}
Thus the centered generalized Gaussian fields $\Phi_{m_0}(x)$ and $\Phi_{M_*}(sx)$ have
identical joint covariances after smearing. This also fixes Wick ordering without a hidden
constant. Indeed, for every regularized centered Gaussian variable $X$,
\[
 {:}e^{tX}{:}=\exp\!\left(tX-\tfrac12t^2\E[X^2]\right),
\]
and the coefficient of $t^n$ is ${:}X^n{:}/n!$. Apply the identity to a common mollifier,
transported by $x\mapsto sx$, and then take the defining distributional limit. Equality of the
regularized covariances shows that the unitary Gaussian identification sends
${:}\mathscr P(\Phi_{m_0}(x)){:}_{m_0}$ to
${:}\mathscr P(\Phi_{M_*}(sx)){:}_{M_*}$ with the coefficients of $\mathscr P$ unchanged.
For a Euclidean cutoff $h$ the change of variables $y=sx$ therefore gives
\[
 \lambda\int_{\R^2}h(x){:}\mathscr P(\Phi_{m_0}(x)){:}_{m_0}\,dx
 =s^{-2}\lambda\int_{\R^2}h(y/s){:}\mathscr P(\Phi_{M_*}(y)){:}_{M_*}\,dy .
\]
The transformed cutoff $h_s(y):=h(y/s)$ still satisfies $0\le h_s\le1$ and has compact
support. Hence the transformed parameters are exactly those displayed in the proposition,
as stated in \cite[p.~590]{GJS}. Finally,
\[
 s^{-2}\lambda=\frac{M_*^2}{m_0^2}\lambda
 <M_*^2\varepsilon_{\mathscr P}=\lambda_*.
\]
This is precisely the source smallness condition. If the source decay rate in the scaled
coordinate $y=sx$ is $m_{1,*}>0$, then
$e^{-m_{1,*}d_y}=e^{-sm_{1,*}d_x}$; thus the rate in the original coordinates is
$m_1:=sm_{1,*}>0$. The transported finite prefactors are the constants named above.
\end{proof}

\subsection{Facts of the cutoff construction}\label{subsec:facts}

\begin{srcfact}\label{fact:sa}
$K(g)$ is essentially self-adjoint and bounded below, and $E(g)$ is a \emph{simple} eigenvalue,
for every measurable $0\le g\le1$ of compact support; $\Om_g$ may be chosen $>0$ a.e.\ in
$Q$-space. \cite[p.~588]{GJS}; \cite{GJ1,GJ2,Rosen70,GJ71}
\end{srcfact}

\begin{srcfact}[dynamics; the side conditions (P), (T)]\label{fact:PT}
There is a one-parameter automorphism group $\alpha$ of $\A$, independent of $g$, with
\begin{align}
  &\alpha_t(\A(O))\subseteq\A(O_{|t|}),\qquad
  \alpha_t(A)=\beta^g_t(A)\ \ \text{if } g\equiv1 \text{ on } O_{|t|};
  \tag{P}\label{eq:P}\\
  &\omega_{\tau_ag}=\omega_g\circ\tau_{-a},
  \tag{T}\label{eq:T}
\end{align}
$\tau_a$ the spatial translation. \eqref{eq:P} is the Glimm--Jaffe patching construction
\cite[\S III]{GJ1}, \cite[\S3.5--3.6]{GJ2}, \cite[\S4.1]{GJExp}; \eqref{eq:T} follows from
$K(\tau_ag)=U_aK(g)U_a^*$ and simplicity (companion M0, Lem.~3.2 --- three lines).
\end{srcfact}

\begin{proposition}[domain smoothing; cutoff propagation]\label{prop:standard-main}
\emph{(i)} $\Om_g\in D(V(h))$ for every $g\in\G$ and real $h\in C_c^\infty$.
\emph{(ii)} $\beta^g_t(\A(O))\subseteq\A(O_{|t|})$ for every $g$, with speed $\le1$ independent
of $g$.
\end{proposition}

\begin{proof}
(i) The model-specific smoothing theorem \cite[Ch.~7, Thm.~7.4]{GJExp} states that, for every
$t>0$, the range of the cutoff semigroup is contained in the domain of each relevant polynomial
multiplication operator; in the present notation,
$\mathrm{Ran}\,e^{-tK(g)}\subseteq D(V(h))$. (Replacing $K(g)$ by $H(g)=K(g)-E(g)$ multiplies
the semigroup by the nonzero scalar $e^{tE(g)}$ and does not change its range.) Since
$K(g)\Om_g=E(g)\Om_g$,
\[
  \Om_g=e^{tE(g)}e^{-tK(g)}\Om_g\in\mathrm{Ran}\,e^{-tK(g)}\subseteq D(V(h)).
\]
(ii) The interaction automorphism $\mathrm{Ad}\,e^{itV(g)}$ has speed zero: for $A$ generated by
Weyl operators supported in $B\Subset O$ pick $\chi\in C_c^\infty(O)$, $\chi=1$ near $B$; then
$V(g)=V(\chi g)+V((1-\chi)g)$ strongly commuting, $e^{itV(\chi g)}\in\A(O)$,
$e^{itV((1-\chi)g)}\in\A(B)'$, so $\mathrm{Ad}\,e^{itV(g)}(A)\in\A(O)$; extend by strong density.
The free dynamics has speed one (Klein--Gordon). The $n$-th Trotter approximant of $\beta^g_t(A)$
then lies in $\A(O_{|t|})$, which is strongly closed, and the Trotter formula converges
strongly. \cite[Ch.~4]{GJExp}
\end{proof}

\section{The source theorems and their cutoff-uniformity}\label{sec:sources}

The Euclidean input is \cite{GJS}. Let $d\Phi$ be the Gaussian with covariance
$(-\Delta_{\R^2}+m_0^2)^{-1}$ and, for $h$ in the class
$\mathcal C=\{0\le h\le1,\ \supp h\ \text{compact}\}$ \cite[p.~589]{GJS},
\[
  dq_h=Z(h)^{-1}e^{-V(h)}d\Phi,\qquad V(h)=\lambda\int{:}\mathscr P(\Phi(z)){:}\,h(z)\,dz .
\]
The observable class (1.1.5) consists of sharp-time monomials
$Q(x^0)=\int d\vec x\,w(\vec x)\prod_i{:}\Phi(x_i)^{m_i}{:}$ with $L^2$ kernels, equal times
allowed, normed by (1.1.8)--(1.1.9) (kernel $L^2$ norms with local-degree weights
$K_2^{2N(\Delta)}N(\Delta)!$ and a single overall prefactor $K_1$; printed forms
verified against the page image of p.~594:
$\|Q\|=K_1[\prod_\Delta K_2^{2N(\Delta)}N(\Delta)!]\,\|w\|_2$ and likewise at fixed times).
Replacing $K_2$ by $\max\{1,K_2\}$ only enlarges the right sides of all source estimates, so
we fix, from now on, $K_2\ge1$.

\begin{srcfact}[CE2 = {\cite[Thm.~1.1.7]{GJS}}, p.~594]\label{sf:CE2}
For localized monomials $Q_1,Q_2$ with disjoint localization squares separated by a strip of
width $d$ bounded by two parallel lines (any orientation),
\[
  \Bigl|\int Q_1Q_2\,dq_h-\int Q_1dq_h\int Q_2dq_h\Bigr|\le C_7e^{-m_1d}\|Q_1\|\|Q_2\| ,
\]
with $m_1,C_7$ \emph{independent of $h$, $Q_1$, $Q_2$} (so stated).
\end{srcfact}

\begin{srcfact}[CE1 = {\cite[Thm.~1.1.8]{GJS}}, p.~595]\label{sf:CE1}
$\bigl|\int Q(x^0)\,dq_h\bigr|\le C_8\|Q(x^0)\|$, uniformly in $h\in\mathcal C$.
\end{srcfact}

\begin{remark}[the uniformity audit]\label{rem:audit}
The class-uniformity of Source Facts \ref{sf:CE2}--\ref{sf:CE1} over all of $\mathcal C$ ---
not merely along $h\uparrow1$ --- is (a) \emph{printed}: \cite[p.~629]{GJS} states ``the
estimates of Theorems 1.1.8, 1.1.11, 3.1 and 4.1 are uniform in the space cutoff $h$'';
(b) \emph{structural}: the convergence proof is a kernel-norm argument in which $h$ enters
quantitatively only through vertex factors majorised by $|\lambda h(x)|\le\lambda$ at first
contact, while every $e^{-V(h)}$ cancels against the normalisation in the final division
(their (4.3)--(4.4)), so no stability estimate enters any constant; and (c) the only
$h$-sensitive analytic prerequisites --- $e^{-\lambda|a_0||K|}\le Z(h)<\infty$ and
$\sup_{h\in\mathcal C,\supp h\subseteq K}\|e^{-V(h)}\|_{L^p(d\Phi)}<\infty$ --- are proved,
self-contained, in the companion audit (\texttt{m4/r1-check.tex}, Lem.~4.2, Thm.~4.5), the two
halves of $0\le h\le1$ entering complementarily (positivity in the Wick lower bound
${:}\mathscr P{:}_c\ge-b(1+c)^{p'}$, boundedness in the UV kernel estimates). The formerly
illegible prefactor of the norm (Gate R2) has been settled by direct inspection of the page
image: $K_1$ enters to the first power and the local weight is $K_2^{2N(\Delta)}$; see
\texttt{m4/r2-check.tex}. All constants below use this exact reading.
\end{remark}

\begin{remark}[conventions lock]\label{rem:conv}
The sharp-time restriction of the Euclidean covariance equals $C$:
$\int\frac{dk_0}{2\pi}(k_0^2+\mu(k)^2)^{-1}=(2\mu(k))^{-1}$ exactly. Hence sharp-time
$d\Phi$-Wick ordering coincides, constant by constant, with the $C$-Wick ordering used in
$V(h)$, and sharp-time functionals of $\Phi(0,\cdot)$ are multiplication operators on $\HH$. No
Wick-mass conversion arises. (Companion \texttt{m4/t3.tex}, Lem.~2.1; \texttt{m2-audit.tex},
Prop.~2.7.)
\end{remark}

Throughout \S\S\ref{sec:M2}--\ref{sec:G}, $g\in\G$ is fixed where limits $T\to\infty$ are taken;
\emph{no uniformity in $g$ beyond the printed constants is used, and no gap is assumed}. Set
\[
  h_{T,g}:=\mathbf 1_{[-T,T]}\otimes g\in\mathcal C,\qquad
  \psi_s:=e^{-sH(g)}\Om_0,\quad\hat\psi_s:=\psi_s/\|\psi_s\|,\quad
  c_g:=\ip{\Om_g}{\Om_0}>0
\]
($c_g>0$ by Source Fact~\ref{fact:sa}: $\Om_g>0$ a.e., $\Om_0\equiv1$).

\begin{lemma}[slab FKN]\label{lem:FKN-main}
For $T>0$, $t\in[0,2T]$, $s\in[-T,T]$, and nonnegative measurable functions $F_1,F_2$ of the
time-zero field:
\begin{align}
  \int F_1\bigl(\Phi(-\tfrac t2,\cdot)\bigr)F_2\bigl(\Phi(\tfrac t2,\cdot)\bigr)dq_{h_{T,g}}
  &=\frac{\ip{F_1\psi_{T-t/2}}{e^{-tH(g)}F_2\psi_{T-t/2}}}{\|\psi_T\|^2},\label{eq:fkn2}\\
  \int F_1\bigl(\Phi(s,\cdot)\bigr)\,dq_{h_{T,g}}
  &=\frac{\ip{\psi_{T-s}}{F_1\,\psi_{T+s}}}{\|\psi_T\|^2},\label{eq:fkn1}
\end{align}
as identities in $[0,\infty]$; energy shifts cancel in the ratios.
\end{lemma}

\begin{proof}
For bounded $F_i$ this is the Feynman--Kac--Nelson formula for
$Z^{-1}e^{-V(h_{T,g})}d\Phi$ (Markov property of $d\Phi$ along constant-time lines;
\cite[Thm.~III.6, Cor.~V.14, Thm.~V.15]{SimonBook};
\cite[Thms.~II.14, II.16]{GRS75}; the identification of sharp-time Wick functionals with
$Q$-space multiplications is Remark~\ref{rem:conv}). For bounded nonnegative $F_i$ both sides
are finite. For bounded signed real $F_i$, write $F_i=F_i^+-F_i^-$ and expand both sides into
the same four finite nonnegative terms; bilinearity proves the identities.

For nonnegative measurable $F_i$, apply the bounded formula to $F_i\wedge M$ and let
$M\uparrow\infty$; monotone convergence on both sides gives the displayed identities in
$[0,\infty]$. For the signed unbounded functions used below, define the clipped truncation
\[
  F^{[M]}:=(-M)\vee(F\wedge M).
\]
Whenever $F\psi_a\in\HH$, Hilbert-space convergence follows from
\[
 \|(F-F^{[M]})\psi_a\|^2
 =\int |F-F^{[M]}|^2|\psi_a|^2\,d\phi_0\longrightarrow0,
 \qquad |F-F^{[M]}|^2|\psi_a|^2\le |F|^2|\psi_a|^2.
\]
For two path-space factors the pointwise estimate
\[
 |F_1(\Phi(s_1))F_2(\Phi(s_2))|
 \le\tfrac12\bigl(F_1(\Phi(s_1))^2+F_2(\Phi(s_2))^2\bigr)
\]
provides a dominated-convergence majorant whenever the two square moments are finite. Source
Fact~\ref{sf:CE1} supplies precisely those square moments for every later application. Thus the
signed identities used below are limits of the bounded signed identities, with no formal
manipulation of products of unbounded operators.
\end{proof}

\begin{lemma}[slab vacuum; simplicity only]\label{lem:vac-main}
$\hat\psi_s\to\Om_g$ in norm, $\|\psi_s\|\downarrow c_g$, and
$\|\psi_{s'}\|/\|\psi_s\|\in[1,1/c_g]$ for $0\le s'\le s$; the ratio tends to $1$ whenever
$s,s'\to\infty$.
\end{lemma}

\begin{proof}
Set $P_g=|\Om_g\rangle\langle\Om_g|$. Since
$P_g\Om_0=c_g\Om_g$,
\[
 \psi_s=c_g\Om_g+e^{-sH(g)}(1-P_g)\Om_0,
\]
and the two summands are orthogonal. If $\mu_g^\perp$ is the spectral measure of
$(1-P_g)\Om_0$, then
\[
 \|e^{-sH(g)}(1-P_g)\Om_0\|^2
 =\int_{(0,\infty)}e^{-2sE}\,d\mu_g^\perp(E)\longrightarrow0
\]
by dominated convergence. The measure has no atom at $0$ by simplicity alone
(Source Fact~\ref{fact:sa}); no isolation and hence no gap is used. It follows that
$\psi_s\to c_g\Om_g$, $\|\psi_s\|\to c_g$, and $\hat\psi_s\to\Om_g$. The spectral formula also
shows that $s\mapsto\|\psi_s\|$ is decreasing. Since
$c_g\le\|\psi_s\|\le\|\psi_0\|=1$, for $0\le s'\le s$,
\[
 1\le\frac{\|\psi_{s'}\|}{\|\psi_s\|}\le\frac1{c_g}.
\]
If both indices tend to infinity, numerator and denominator tend to $c_g>0$, so their ratio
tends to $1$.
\end{proof}

\section{The uniform local interaction moment}\label{sec:M2}

\begin{theorem}[(M$_2$)]\label{thm:M2-main}
There is $M=M(\lambda,\mathscr P,K_1,K_2,C_8)<\infty$ such that for every $g\in\G$ and every
real $h\in C_c^\infty(\R)$ with $\|h\|_\infty\le1$ and $\supp h$ contained in an interval of
length $2$,
\[
  \|V(h)\Om_g\|\ \le\ M .
\]
Explicitly $M^2=9\,C_8\,\mathfrak K(p')\,\lambda^2(\sum_n|a_n|)^2$ with
$\mathfrak K(p')=K_1K_2^{8p'}(4p')!$ (per-point degree $\le\deg\mathscr P=2p'$, so the local degrees sum to $n+n'\le4p'$).
\end{theorem}

\begin{proof}
\emph{Step 1 (class membership and norm).} On path space set
\[
  Q_V:=\lambda^2\!\!\sum_{n,n'=0}^{2p'}\!\!a_na_{n'}\iint h(x)h(y)\,
  {:}\Phi(0,x)^n{:}\,{:}\Phi(0,y)^{n'}{:}\,dx\,dy ,
\]
a sum of equal-time monomials of the class (1.1.5) with $r=2$ and kernels
$a_na_{n'}h\otimes h\in L^2$ --- \emph{no Wick recombination occurs}, the product of two
individually ordered powers being literally of the form (1.1.5); per-point degrees are
$\le2p'$, admissible by \cite[fn.~9]{GJS}. Let
$J:=\{j\in\Z:\|h\mathbf1_{[j,j+1)}\|_2>0\}$ and decompose
$h=\sum_{j\in J}h\mathbf1_{[j,j+1)}$. An interval of length $2$ meets at most three half-open
unit cells in positive measure, so $|J|\le3$, while
$\|h\mathbf1_{[j,j+1)}\|_2\le\|h\|_\infty\le1$. A localized piece with both points in one square carries
$K_2^{2(n+n')}(n+n')!\le K_2^{8p'}(4p')!$; with points in distinct squares,
$K_2^{2n}n!\,K_2^{2n'}n'!\le K_2^{8p'}(4p')!$ --- the $K_2$ exponent $2\sum_\Delta N(\Delta)
=2(n+n')\le8p'$ is additive and case-independent, and $(n+n')!\le(4p')!$ since each degree is
at most $\deg\mathscr P=2p'$. Summing the $\le9$ pairs and the degrees,
$\|Q_V\|\le9\mathfrak K(p')\lambda^2(\sum|a_n|)^2$.

\emph{Step 2 (uniform slab bound).} Apply \eqref{eq:fkn2} at $t=0$ first to
$F_1=F_2=|V(h)|\wedge R$. Both functions are nonnegative. Monotone convergence as
$R\uparrow\infty$, followed by Source Fact~\ref{sf:CE1}, gives
(Remark~\ref{rem:conv} identifies the sharp-time functional with the operator)
\[
  \|V(h)\hat\psi_T\|^2=\int Q_V\,dq_{h_{T,g}}\ \le\ C_8\|Q_V\|\ \le\ M^2
  \qquad\text{for all }T>0,\ g\in\G ,
\]
since $h_{T,g}\in\mathcal C$.

\emph{Step 3 (closedness limit).} $V(h)$ is multiplication by a real measurable function,
hence self-adjoint on its maximal domain and closed. Its graph is a norm-closed linear subspace
of $\HH\oplus\HH$, hence is weakly closed. Choose $T_k\to\infty$ so that
$\|V(h)\hat\psi_{T_k}\|$ tends to its liminf. Step 2 gives a weakly convergent subsequence,
say $V(h)\hat\psi_{T_{k_\ell}}\rightharpoonup y$, while
Lemma~\ref{lem:vac-main} gives $\hat\psi_{T_{k_\ell}}\to\Om_g$ strongly. Weak closedness of the
graph yields $\Om_g\in D(V(h))$ and $y=V(h)\Om_g$. Therefore weak lower semicontinuity gives
\[
 \|V(h)\Om_g\|\le\liminf_{\ell\to\infty}
 \|V(h)\hat\psi_{T_{k_\ell}}\|
 =\liminf_{T\to\infty}\|V(h)\hat\psi_T\|\le M.
\]
\end{proof}

\section{The uniform gap}\label{sec:G}

Let $\Q$ be the set of finite products $Q=\prod_{i=1}^r\varphi_0(f_i)$, $f_i\in C_c^\infty$
real, viewed simultaneously as multiplication operators on $\HH$ and, at Euclidean time $s$, as
class-(1.1.5) monomials $Q(s)$ with kernel $f_1\otimes\cdots\otimes f_r$ and all $m_i=1$. Write
$\|Q\|_{\rm GJS}$ for the (1.1.9)-norm of $Q(s)$. Let $I$ be a bounded interval of length
$L$ containing every $\supp f_i$, and define
\[
 J_I:=\{j\in\Z:|I\cap[j,j+1)|>0\},\qquad N_I:=\#J_I .
\]
If $j_{\min}$ and $j_{\max}$ are the extreme elements of $J_I$, then points of positive-measure
intersection can be chosen arbitrarily close to the two extreme cells' interiors; consequently
$j_{\max}-j_{\min}-1<L$. Since $N_I=j_{\max}-j_{\min}+1$ is an integer,
\[
  N_I\le\lfloor L\rfloor+2\le L+2.
\]
Decompose $f_i=\sum_{j\in J_I}f_i\mathbf1_{[j,j+1)}$. For every cell assignment, the product
of the local factorials is at most $r!$, because
$\prod_\Delta N(\Delta)!\le(\sum_\Delta N(\Delta))!=r!$, and the $K_2$ exponent is $2r$.
Moreover Cauchy--Schwarz gives
\[
 \sum_{j\in J_I}\|f_i\mathbf1_{[j,j+1)}\|_2
 \le N_I^{1/2}\Bigl(\sum_{j\in J_I}
       \|f_i\mathbf1_{[j,j+1)}\|_2^2\Bigr)^{1/2}
 =N_I^{1/2}\|f_i\|_2.
\]
Summing all cell assignments, and applying the same calculation to the $2r$ factors of $Q^2$,
yields the exact bounds
\begin{equation}\label{eq:normbound}
  \|Q\|_{\rm GJS}\le K_1K_2^{2r}r!\,N_I^{r/2}\prod_{i=1}^r\|f_i\|_2,\qquad
  \|Q^2\|_{\rm GJS}\le K_1K_2^{4r}(2r)!\,N_I^r\prod_i\|f_i\|_2^2
  =:\mathfrak n(Q)^2 .
\end{equation}
This norm is independent of the common sharp time $s$: at every $s$ all factors occupy one
temporal row, and changing $s$ changes only that common row label. The spatial cell assignments,
local degrees, kernel norms, and local weights in the preceding calculation are unchanged; no
invariance of a fixed unit lattice under arbitrary continuous translation is asserted.

For $t>2$, the pieces of $Q(-t/2)$ and of $\tilde Q(t/2)$ are separated by a horizontal strip of
width $\ge t-2$ (unit lattice rows), so Source Fact~\ref{sf:CE2} and bilinearity give the
\emph{two-time cluster bound}
\begin{equation}\label{eq:2time}
  \Bigl|\int Q(-\tfrac t2)\tilde Q(\tfrac t2)dq_h-\int Q(-\tfrac t2)dq_h
  \int\tilde Q(\tfrac t2)dq_h\Bigr|
  \le C_7e^{2m_1}e^{-m_1t}\|Q\|_{\rm GJS}\|\tilde Q\|_{\rm GJS},\quad h\in\mathcal C .
\end{equation}

\begin{lemma}\label{lem:Gaux}
Fix $g\in\G$ and $Q\in\Q$. Then:
\begin{enumerate}
\item[\rm(i)] $\sup_T\|Q\hat\psi_T\|^2\le C_8\,\mathfrak n(Q)^2$;
      $\Om_g\in D(Q)$ and $\|Q\Om_g\|^2\le C_8\mathfrak n(Q)^2$;
      for every $a>0$, the clipped truncations
      $Q^{[M]}:=(-M)\vee(Q\wedge M)$ satisfy
      $\|(Q-Q^{[M]})\psi_a\|\to0$;
\item[\rm(ii)] for every fixed $s\in\R$,
      $\int Q(s)\,dq_{h_{T,g}}\to\omega_g(Q)$ as $T\to\infty$ through $T>|s|$;
\item[\rm(iii)] for every $t\ge0$,
      $\displaystyle
      \ip{Q\Om_g}{e^{-tH(g)}Q\Om_g}\le\liminf_{T\to\infty}\int Q(-\tfrac t2)Q(\tfrac t2)
      \,dq_{h_{T,g}}$.
\end{enumerate}
\end{lemma}

\begin{proof}
\emph{(i).} Apply \eqref{eq:fkn2} at $t=0$ to
$F_1=F_2=|Q|\wedge M$. Monotone convergence and Source Fact~\ref{sf:CE1}, applied to the
class-(1.1.5) observable $Q^2$, give
\begin{equation}\label{eq:Gaux-bound}
 \|Q\hat\psi_T\|^2
 =\int Q(0)^2\,dq_{h_{T,g}}
 \le C_8\|Q^2\|_{\rm GJS}
 =C_8\mathfrak n(Q)^2.
\end{equation}
Choose a sequence attaining the liminf as $T\to\infty$ and repeat the weak-graph argument of
Theorem~\ref{thm:M2-main}, Step 3; this proves $\Om_g\in D(Q)$ and, by weak lower
semicontinuity, $\|Q\Om_g\|^2\le C_8\mathfrak n(Q)^2$. For fixed $a>0$,
\eqref{eq:Gaux-bound} with
$T=a$ gives $Q\psi_a\in\HH$, and dominated convergence on $Q$-space gives
\[
 \|(Q-Q^{[M]})\psi_a\|^2
 =\int |Q-Q^{[M]}|^2|\psi_a|^2\,d\phi_0\longrightarrow0,
 \quad |Q-Q^{[M]}|^2|\psi_a|^2\le Q^2|\psi_a|^2.
\]

\emph{(ii).} Apply the bounded signed identity \eqref{eq:fkn1} to $Q^{[M]}$ and pass
$M\to\infty$ using part (i) on the Hilbert-space side. On path space, the positive and negative
parts are integrable because
$\int|Q(s)|\,dq\le(\int Q(s)^2\,dq)^{1/2}$ and Source Fact~\ref{sf:CE1} bounds the square.
Thus, for $T>|s|$,
\begin{equation}\label{eq:onepoint-transfer}
 \int Q(s)\,dq_{h_{T,g}}
 =\Bigl\langle\hat\psi_{T-s},
      \frac{Q\psi_{T+s}}{\|\psi_T\|}\Bigr\rangle
      \frac{\|\psi_{T-s}\|}{\|\psi_T\|}.
\end{equation}
The first vector tends strongly to $\Om_g$, and the final scalar tends to $1$ by
Lemma~\ref{lem:vac-main}. For the middle vector, write
\begin{equation}\label{eq:Gaux-middle}
 \frac{Q\psi_{T+s}}{\|\psi_T\|}
 =\frac{\|\psi_{T+s}\|}{\|\psi_T\|}\,Q\hat\psi_{T+s}.
\end{equation}
The scalar ratio in \eqref{eq:Gaux-middle} is at most $1$ when $s\ge0$. When $s<0$,
$T+s<T$ and Lemma~\ref{lem:vac-main} gives instead
\[
 1\le\frac{\|\psi_{T+s}\|}{\|\psi_T\|}\le c_g^{-1}.
\]
Hence \eqref{eq:Gaux-bound} makes the middle vectors bounded in both sign cases. The scalar ratio tends to
$1$ for fixed $s$. Every weakly convergent subsequence of
$Q\hat\psi_{T+s}$ has limit $Q\Om_g$: indeed,
$\hat\psi_{T+s}\to\Om_g$ strongly and the graph of $Q$ is weakly closed. Thus
$Q\hat\psi_{T+s}\rightharpoonup Q\Om_g$, and \eqref{eq:Gaux-middle} has the same weak limit. Taking the
strong--weak pairing in \eqref{eq:onepoint-transfer} proves
$\int Q(s)\,dq_{h_{T,g}}\to\ip{\Om_g}{Q\Om_g}$.

\emph{(iii).} Apply \eqref{eq:fkn2} first to the bounded signed functions $Q^{[M]}$ at the
two time slices. The path-space products converge in $L^1$ because
\[
 |Q^{[M]}(-t/2)Q^{[M]}(t/2)|
 \le\tfrac12\bigl(Q(-t/2)^2+Q(t/2)^2\bigr),
\]
whose integral is finite by Source Fact~\ref{sf:CE1}. Part (i) gives convergence of the
Hilbert-space vectors. Consequently, for $T>t/2$,
\begin{equation}\label{eq:Gaux-two}
 \int Q(-t/2)Q(t/2)\,dq_{h_{T,g}}
 =\|e^{-tH(g)/2}u_T\|^2,\qquad
 u_T:=\frac{Q\psi_{T-t/2}}{\|\psi_T\|}.
\end{equation}
Here
\[
 \|u_T\|
 =\frac{\|\psi_{T-t/2}\|}{\|\psi_T\|}\,
   \|Q\hat\psi_{T-t/2}\|
 \le c_g^{-1}C_8^{1/2}\mathfrak n(Q)
\]
by Lemma~\ref{lem:vac-main} and \eqref{eq:Gaux-bound}. The ratio tends to $1$, and the same weak-graph
argument as in (ii) gives $u_T\rightharpoonup Q\Om_g$. Since $e^{-tH(g)/2}$ is bounded,
$e^{-tH(g)/2}u_T\rightharpoonup e^{-tH(g)/2}Q\Om_g$. Weak lower semicontinuity of the norm
in \eqref{eq:Gaux-two} now yields
\[
 \ip{Q\Om_g}{e^{-tH(g)}Q\Om_g}
 =\|e^{-tH(g)/2}Q\Om_g\|^2
 \le\liminf_{T\to\infty}\int Q(-t/2)Q(t/2)\,dq_{h_{T,g}}.
\]
\end{proof}

\begin{proposition}[connected decay]\label{prop:decay-main}
For real $Q\in\Q$ and $\psi_Q:=Q\Om_g-\omega_g(Q)\Om_g$:
\[
  \ip{\psi_Q}{e^{-tH(g)}\psi_Q}\le\widehat C_Qe^{-m_1t}\quad(t\ge0),\qquad
  \widehat C_Q=\max\bigl(C_7e^{2m_1}\|Q\|_{\rm GJS}^2,\ e^{3m_1}\|\psi_Q\|^2\bigr).
\]
\end{proposition}

\begin{proof}
$\ip{\psi_Q}{e^{-tH}\psi_Q}=\ip{Q\Om_g}{e^{-tH}Q\Om_g}-\omega_g(Q)^2$ (using
$e^{-tH}\Om_g=\Om_g$ and $Q$ symmetric). For $t>2$, Lemma~\ref{lem:Gaux}(iii), then
\eqref{eq:2time} with $h=h_{T,g}$ uniformly in $T$, then Lemma~\ref{lem:Gaux}(ii) for the
disconnected part, give
\[
  \ip{Q\Om_g}{e^{-tH}Q\Om_g}\le C_7e^{2m_1}e^{-m_1t}\|Q\|_{\rm GJS}^2+\omega_g(Q)^2 ,
\]
via $\liminf(a_T+b_T)\le\sup_Ta_T+\lim_Tb_T$. For $t\le3$ use monotonicity of
$t\mapsto\|e^{-tH/2}\psi_Q\|^2$ and $e^{-m_1t}\ge e^{-3m_1}$.
\end{proof}

\begin{theorem}[density]\label{thm:dense-main}
$\Span_{\mathbb C}\{Q\Om_g:Q\in\Q\}$ is dense in $\HH$, for every $g\in\G$.
\end{theorem}

\begin{proof}
Equivalently (unitarity of $\psi\mapsto\psi/\Om_g$ onto $L^2(d\mu_g)$, $d\mu_g=\Om_g^2d\phi_0$):
polynomials in $\{\varphi_0(f)\}$ are dense in $L^2(d\mu_g)$.

\emph{Exponential moments.} Let $f$ be real and supported in an interval $I$; use the integer
$N_I$ defined above and set
\[
 X_f:=\varphi_0(f),\qquad K_f:=K_2^2N_I^{1/2}\|f\|_2.
\]
If $f=0$, then $X_f=0$ and all exponential moments are finite; hence assume $f\ne0$, so
$K_f>0$.
For every $n\ge1$, Lemma~\ref{lem:Gaux}(i), applied to $Q=X_f^n$, and
\eqref{eq:normbound} give
\[
 \E_{\mu_g}[X_f^{2n}]
 =\|X_f^n\Om_g\|^2
 \le C_8K_1K_2^{4n}(2n)!\,N_I^n\|f\|_2^{2n}
 =C_8K_1K_f^{2n}(2n)!.
\]
Since $e^{s|x|}\le e^{sx}+e^{-sx}=2\cosh(sx)$, Tonelli's theorem and
$\E_{\mu_g}1=1$ yield, for $0\le s<K_f^{-1}$,
\[
 \E_{\mu_g}e^{s|X_f|}
 \le2\left(1+C_8K_1\sum_{n=1}^{\infty}(sK_f)^{2n}\right)
 \le2\max\{1,C_8K_1\}\sum_{n=0}^{\infty}(sK_f)^{2n}<\infty.
\]
Thus $K_1$ affects the prefactor but not the nonzero radius of exponential integrability.

\emph{Orthogonality kills conditional expectations.} Suppose
$\zeta\in L^2(d\mu_g)$ is orthogonal to every cylinder polynomial. Fix real
$e_1,\dots,e_k\in C_c^\infty$, put $X_j=\varphi_0(e_j)$, and choose $r_j>0$ with
$\E_{\mu_g}e^{r_j|X_j|}<\infty$ by the preceding paragraph. Choose $\sigma>0$ so that
$2k\sigma<\min_j r_j$. For $|\operatorname{Re}z_j|<\sigma$ define
\[
 F(z_1,\dots,z_k)
 :=\E_{\mu_g}\!\left[\overline\zeta\,
       \exp\!\left(\sum_{j=1}^kz_jX_j\right)\right].
\]
Cauchy--Schwarz followed by H\"older with $k$ equal exponents gives the explicit dominator
\begin{align*}
 \E_{\mu_g}\!\left[|\zeta|e^{\sum_j\operatorname{Re}z_jX_j}\right]
 &\le\|\zeta\|_2
       \left(\E_{\mu_g}e^{2\sigma\sum_j|X_j|}\right)^{1/2}\\
 &\le\|\zeta\|_2
       \prod_{j=1}^k\left(\E_{\mu_g}e^{2k\sigma|X_j|}\right)^{1/(2k)}
 <\infty.
\end{align*}
On every closed smaller polystrip, derivatives are dominated by the same expression with a
slightly larger exponent, because for $a>0$ and $n\in\mathbb N$,
$|x|^n\le n!a^{-n}e^{a|x|}$. Hence differentiation under the integral is justified and $F$ is
jointly analytic. Every derivative at the origin equals the inner product of $\zeta$ with a
cylinder monomial and therefore vanishes. Applying the one-variable identity theorem
successively in $z_1,\ldots,z_k$ gives $F\equiv0$ on the connected polystrip, in particular
on $(i\R)^k$.

Define the finite complex Borel measure on $\R^k$
\[
 \nu(B):=\E_{\mu_g}\!\left[\overline\zeta\,
   \mathbf1_{\{(X_1,\ldots,X_k)\in B\}}\right].
\]
Its Fourier transform is $F(it_1,\ldots,it_k)=0$. Uniqueness of Fourier transforms of finite
measures gives $\nu=0$, which is equivalent to
$\E_{\mu_g}[\zeta\mid\Sigma_k]=0$ for
$\Sigma_k:=\sigma(X_1,\ldots,X_k)$.

\emph{Martingale exhaustion.} Choose a sequence
$(e_j)\subset C_c^\infty(\R;\R)$ dense in real $L^2(\R)$ and let
$\Sigma_k=\sigma(\varphi_0(e_1),\ldots,\varphi_0(e_k))$. The bounded free covariance
$C=(2\sqrt{-\partial_x^2+m_0^2})^{-1}$ gives
\[
 \E_{\phi_0}|\varphi_0(f)-\varphi_0(e)|^2
 =\ip{f-e}{C(f-e)}
 \le\|C\|\,\|f-e\|_2^2.
\]
For any $f\in L^2(\R;\R)$, an approximating sequence has a subsequence converging
$\phi_0$-almost surely and therefore $\mu_g$-almost surely, because
$d\mu_g=\Om_g^2d\phi_0\ll d\phi_0$. Thus every cylinder coordinate is measurable modulo
$\mu_g$-null sets with respect to $\Sigma_\infty:=\bigvee_k\Sigma_k$, so
$\Sigma_\infty$ is the full cylinder sigma-algebra. The $L^2$ martingale convergence theorem
now gives
\[
 \zeta=\E_{\mu_g}[\zeta\mid\Sigma_\infty]
 =\lim_{k\to\infty}\E_{\mu_g}[\zeta\mid\Sigma_k]=0.
\]
\end{proof}

\begin{theorem}[(G)]\label{thm:G-main}
For every $g\in\G$:\quad
$\spec H(g)\subseteq\{0\}\cup[m_1,\infty)$, and $0$ is simple.
\end{theorem}

\begin{proof}
Let $D=\{\psi_Q:Q\in\Q\text{ real}\}$. The complex span of
$D\cup\{\Om_g\}$ equals the complex span of $\{Q\Om_g:Q\in\Q\}$ and is dense by
Theorem~\ref{thm:dense-main}. Fix $0\le\beta<m_1$ and write
$P_\beta:=\mathbf1_{[0,\beta]}(H(g))$. For $\psi_Q\in D$, the spectral theorem and
Proposition~\ref{prop:decay-main} give, for every $t\ge0$,
\[
 e^{-\beta t}\|P_\beta\psi_Q\|^2
 \le\ip{\psi_Q}{e^{-tH(g)}\psi_Q}
 \le\widehat C_Qe^{-m_1t}.
\]
Multiplication by $e^{\beta t}$ and passage to $t\to\infty$ yield
$P_\beta\psi_Q=0$, because $m_1-\beta>0$. On the other hand,
$P_\beta\Om_g=\Om_g$ because $0\in[0,\beta]$. Thus $P_\beta$ maps a dense subspace into
$\mathbb C\Om_g$. Continuity of the bounded projection extends this containment to all of
$\HH$, and $P_\beta\Om_g=\Om_g$ shows
\[
 P_\beta=|\Om_g\rangle\langle\Om_g|\qquad(0\le\beta<m_1).
\]
If $\spec H(g)$ met $(0,m_1)$, an open interval
$J\Subset(0,m_1)$ meeting the spectrum would have a nonzero spectral projection. Choosing
$\beta<m_1$ above $\sup J$ would give
\[
 0\ne\mathbf1_J(H(g))
 \le P_\beta-|\Om_g\rangle\langle\Om_g|=0,
\]
a contradiction. Hence $\spec H(g)\subseteq\{0\}\cup[m_1,\infty)$; simplicity of the zero
eigenvalue is Source Fact~\ref{fact:sa}.
\end{proof}

\begin{remark}\label{rem:indep}
The constants $\widehat C_Q$ of Proposition~\ref{prop:decay-main} are in fact $g$-uniform, but
this is never used: the spectral argument needs only the \emph{rate} to be $g$-independent.
Accordingly (M$_2$) and (G) are logically independent consequences of the same source pair.
No gap, and no isolation of $E(g)$, was assumed at any point of
\S\S\ref{sec:sources}--\ref{sec:G}; isolation with the uniform bound $m_1$ is part of the
conclusion of Theorem~\ref{thm:G-main}.
\end{remark}

\section{The Cauchy step (LPPL)}\label{sec:lppl}

Throughout this section, (G) and (M) denote the statements supplied by
Theorems~\ref{thm:G-main} and \ref{thm:M2-main}: $\spec H(g)\subseteq\{0\}\cup[\gamma,\infty)$
for every $g\in\G$ with $\gamma:=m_1$, and $\sup_{g}\|V(h)\Om_g\|\le M$ for real
$h\in C_c^\infty$ with $\|h\|_\infty\le1$ and support in an interval of length $2$.

\subsection{Affine-path regularity from FKN}

\begin{proposition}[path regularity]\label{prop:FKNpath-main}
Let $g,g'\in\G$ be equal to one on a common nonempty interval, $g_s=(1-s)g+sg'$, $v=g'-g$, and
\[
  K_s=K(g_s),\quad E_s=E(g_s),\quad H_s=K_s-E_s,\quad
  P_s=|\Om_s\rangle\langle\Om_s|,\quad Q_s=1-P_s .
\]
Under \textup{(G)}, $s\mapsto P_s$ is norm-continuous on $[0,1]$ and $C^1$ on $(0,1)$; locally
on $(0,1)$ normalized ground states can be chosen $C^1$ in the parallel gauge, and then
\begin{equation}\label{eq:HFmain}
  \dot E_s=\ip{\Om_s}{V(v)\Om_s},\qquad
  \dot\Om_s=-H_s^{-1}Q_s\,V(v)\,\Om_s .
\end{equation}
In particular $s\mapsto\omega_{g_s}(A)=\operatorname{Tr}(P_sA)$ is continuous on $[0,1]$ for
every $A\in B(\HH)$.
\end{proposition}

\begin{proof}
Fix $T>0$. On the free Euclidean path space $(\mathcal S'(\R^2),d\mu_0)$, $d\mu_0=d\Phi$, let
$J_t$ denote the time-$t$ slice embedding, $(J_tf)(\Phi)=f(\Phi(t,\cdot))$ for bounded cylinder
$f$, and set
\[
  \mathcal U_s=\lambda\int_0^T\!\!\int_\R g_s(x)\,{:}\mathscr P(\Phi(t,x)){:}\,dx\,dt,
  \qquad w_s=e^{-\mathcal U_s},\qquad \mathcal U_v=\mathcal U_1-\mathcal U_0 .
\]
The Feynman--Kac--Nelson formula reads
$\ip{f}{e^{-TK_s}h}=\int\overline{J_Tf}\,(J_0h)\,w_s\,d\mu_0$ for bounded time-zero $f,h$
\cite[Cor.~V.14, Thm.~V.15]{SimonBook}; \cite[Thm.~II.16]{GRS75}.

\emph{Time-slice product estimate.} Put $q_T=2/(1-e^{-m_0T})$ and $r_T=q_T/(q_T-1)=
2/(1+e^{-m_0T})$, so $p_T:=2/r_T=1+e^{-m_0T}$ has conjugate $p_T'=1+e^{m_0T}$ and
$(p_T'-1)/(p_T-1)=e^{2m_0T}$ --- exactly the Nelson hypercontractivity condition for the free
Markov semigroup $\mathsf P_T=e^{-TH_0}$ on the time-zero Gaussian $\nu_0$:
$\|\mathsf P_TF\|_{p_T'}\le\|F\|_{p_T}$ \cite[Thms.~I.17, III.17]{SimonBook}. The Markov property
and H\"older then give
\[
  \|(J_Tf)(J_0h)\|_{r_T}^{r_T}
  =\int|h|^{r_T}\,\mathsf P_T(|f|^{r_T})\,d\nu_0
  \le\||h|^{r_T}\|_{p_T}\,\||f|^{r_T}\|_{p_T}
  =\|h\|_2^{r_T}\|f\|_2^{r_T},
\]
using $r_Tp_T=2$; hence
\begin{equation}\label{eq:timeslice-main}
  \|(J_Tf)(J_0h)\|_{r_T}\le\|f\|_2\,\|h\|_2 .
\end{equation}

\emph{Weight regularity.} $\mathcal U_v\in L^p(\mu_0)$ for every $p<\infty$, and since
$g,g'\ge0$ and $\mathscr P$ is semibounded, $e^{-a\mathcal U_0},e^{-a\mathcal U_1}\in L^1(\mu_0)$
for every $a<\infty$ \cite[Thm.~V.7]{SimonBook} (this is the $h$-uniform stability of
Remark~\ref{rem:audit}(c)). By $\mathcal U_s=(1-s)\mathcal U_0+s\mathcal U_1$ and H\"older,
\begin{equation}\label{eq:wmoments}
  \|w_s\|_a^a\le\Bigl(\int e^{-a\mathcal U_0}d\mu_0\Bigr)^{1-s}
  \Bigl(\int e^{-a\mathcal U_1}d\mu_0\Bigr)^{s},
\end{equation}
so $\sup_{0\le s\le1}\|w_s\|_a<\infty$. If $s_n\to s$ then $\mathcal U_{s_n}\to\mathcal U_s$ in
every $L^p$, so $w_{s_n}\to w_s$ in probability; \eqref{eq:wmoments} at a larger exponent makes
$(|w_{s_n}|^a)$ uniformly integrable, and Vitali gives $\|w_{s_n}-w_s\|_a\to0$. Pointwise on
path space,
\[
  \frac{w_{s+h}-w_s}{h}=-\mathcal U_v\int_0^1w_{s+\theta h}\,d\theta ,
\]
so for $0<s<1$, choosing $p,a<\infty$ with $q_T^{-1}=p^{-1}+a^{-1}$,
\[
  \Bigl\|\frac{w_{s+h}-w_s}{h}+\mathcal U_vw_s\Bigr\|_{q_T}
  \le\|\mathcal U_v\|_p\sup_{0\le\theta\le1}\|w_{s+\theta h}-w_s\|_a\to0 :
\]
$s\mapsto w_s$ is continuous $[0,1]\to L^{q_T}$ and $C^1$ $(0,1)\to L^{q_T}$ with continuous
derivative $-\mathcal U_vw_s$.

\emph{Transfer to the semigroup.} Since $1/r_T+1/q_T=1$, FKN and
\eqref{eq:timeslice-main} give
$\|e^{-TK_s}-e^{-TK_r}\|\le\|w_s-w_r\|_{q_T}$ (first on bounded $f,h$, then by $L^2$ density),
and the same estimate for difference quotients: $s\mapsto L_s:=e^{-TK_s}$ is norm-continuous on
$[0,1]$, norm-$C^1$ on $(0,1)$.

\emph{Spectral consequences.} The largest spectral value of $L_s$ is the simple eigenvalue
$\lambda_s=e^{-TE_s}=\|L_s\|$, hence continuous; by (G),
$\spec L_s\setminus\{\lambda_s\}\subseteq[0,e^{-T\gamma}\lambda_s]$, a relative separation
locally bounded below. The Riesz integral over a circle about $\lambda_s$ yields norm continuity
of $P_s$ on $[0,1]$ and $C^1$ on $(0,1)$. Near $s_0\in(0,1)$,
$\widetilde\Om_s=P_s\Om_{s_0}/\|P_s\Om_{s_0}\|$ is $C^1$; a $C^1$ phase puts it in the parallel
gauge $\ip{\Om_s}{\dot\Om_s}=0$; and $E_s=-T^{-1}\log\ip{\Om_s}{L_s\Om_s}$ is $C^1$.

\emph{Identification of the derivatives.} On the common core $\mathscr C=C^\infty(H_0)$,
$K_s\xi=K_{s_0}\xi+(s-s_0)V(v)\xi$ \cite[Thm.~V.12(a)]{SimonBook}. Differentiating
$\ip{K_s\xi}{\Om_s}=E_s\ip{\xi}{\Om_s}$ at fixed $\xi\in\mathscr C$, using
$\Om_s\in D(V(v))$ (Proposition~\ref{prop:standard-main}) and self-adjointness of $V(v)$,
\[
  \ip{H_s\xi}{\dot\Om_s}=\ip{\xi}{(\dot E_s-V(v))\Om_s}\qquad(\xi\in\mathscr C).
\]
$\mathscr C$ is a core for $H_s$, so $\dot\Om_s\in D(H_s)$ and
$H_s\dot\Om_s=(\dot E_s-V(v))\Om_s$; pairing with $\Om_s$ gives the first identity of
\eqref{eq:HFmain}, and the parallel gauge with $\|H_s^{-1}Q_s\|\le\gamma^{-1}$ gives the second.
Finally $|\operatorname{Tr}((P_s-P_r)A)|\le\|P_s-P_r\|_1\|A\|\le2\|P_s-P_r\|\|A\|$ (rank
$\le2$), which yields the endpoint continuity assertion.
\end{proof}

\subsection{The exact spectral filter}

\begin{lemma}[filter]\label{prop:filter-main}
For every $\gamma>0$ there is $W_\gamma\in L^1(\R)$ with
$\widehat W_\gamma(E):=\int W_\gamma(t)e^{itE}dt$ satisfying
\[
  \widehat W_\gamma(0)=0,\qquad \widehat W_\gamma(E)=\frac1E\ \ (|E|\ge\gamma),\qquad
  \int(1+|t|)^q|W_\gamma(t)|\,dt<\infty\ \ (\forall q\ge0).
\]
Consequently, if $H\ge0$, $\spec H\subseteq\{0\}\cup[\gamma,\infty)$, $Q=1-P_{\ker H}$, then
$H^{-1}Q=\int W_\gamma(t)e^{itH}Q\,dt$ as a norm-convergent Bochner integral. Write
$T_\gamma(r):=\int_{|t|>r}|W_\gamma|\,dt\le C_{q,\gamma}(1+r)^{-q}$.
\end{lemma}

\begin{proof}
Pick an even $\chi\in C^\infty(\R;[0,1])$, $\chi=0$ on $[-\gamma/2,\gamma/2]$, $\chi=1$ off
$(-\gamma,\gamma)$, and set $\widehat W_\gamma(E)=\chi(E)/E$ ($=0$ at $E=0$): a smooth odd
function whose $k$-th derivative is integrable for every $k\ge1$ (smooth on the compact
transition region, $O(E^{-k-1})$ outside). Define
$W_\gamma(t)=\frac1{2\pi}\int e^{-itE}\widehat W_\gamma(E)dE$ as an oscillatory integral; by
oddness, $W_\gamma(t)=-\frac i\pi\int_0^\infty\sin(tE)\chi(E)E^{-1}dE$, which is uniformly
bounded for $0<|t|\le1$ (absolute bound on $[0,\gamma]$; substitution $u=|t|E$ and Dirichlet's
test on $[\gamma,\infty)$). For $|t|\ge1$, $k$ integrations by parts give
$|W_\gamma(t)|\le\|\widehat W_\gamma^{(k)}\|_{L^1}/(2\pi|t|^k)$, all boundary terms vanishing;
$k>q+1$ gives the moment bound and $W_\gamma\in L^1$. Fourier inversion in $\mathcal S'$ plus
continuity of the transform of an $L^1$ function give
$\mathcal FW_\gamma=\widehat W_\gamma$ pointwise, and the spectral theorem yields
$\int W_\gamma(t)e^{itH}Q\,dt=\widehat W_\gamma(H)Q=H^{-1}Q$, norm-convergent since
$W_\gamma\in L^1$, $\|e^{itH}Q\|\le1$. The tail bound is Markov's inequality.
\end{proof}

\subsection{Clustering with one unbounded local observable}

\begin{lemma}[spectral separation]\label{lem:spectral-main}
Let
\[
 F(t)=\int_{[\gamma,\infty)}e^{itE}d\rho(E),\qquad
 G(t)=\int_{(-\infty,-\gamma]}e^{itE}d\sigma(E),
\]
with
$|\rho|(\R),|\sigma|(\R)\le M_0$. If $F=G$ on $(-R,R)$ and $\gamma R\ge1$, then
$|F(0)|\le2M_0e^{-\gamma R/2}$.
\end{lemma}

\begin{proof}
For $a,\varepsilon>0$ let $f_{a,\varepsilon}(t)=\frac1{2\pi i}e^{-t^2/2a}(t-i\varepsilon)^{-1}$;
residues and Gaussian convolution give
$\check f_{a,\varepsilon}(E)=\int_0^\infty\sqrt{a/2\pi}\,e^{-a(E-E')^2/2}e^{-\varepsilon E'}dE'
\in[0,1]$, with $\check f_{a,\varepsilon}\to\Phi(\sqrt aE)$ as $\varepsilon\downarrow0$
($\Phi$ = normal distribution function), and
$\int_{|t|\ge R}|f_{a,\varepsilon}|\le\frac a{\pi R^2}e^{-R^2/2a}$. Fubini, $F=G$ on $(-R,R)$,
and the variation bounds give
$|\int\Phi(\sqrt aE)d\rho-\int\Phi(\sqrt aE)d\sigma|\le\frac{2M_0a}{\pi R^2}e^{-R^2/2a}$
after $\varepsilon\downarrow0$. With $1-\Phi(x),\Phi(-x)\le\frac12e^{-x^2/2}$ and the support
conditions, $|F(0)|\le M_0e^{-a\gamma^2/2}+\frac{2M_0a}{\pi R^2}e^{-R^2/2a}$; $a=R/\gamma$ and
$\gamma R\ge1$ give the claim.
\end{proof}

\begin{lemma}[one unbounded argument]\label{prop:unbdd-main}
Let $H\ge0$ have simple normalized ground state $\Om$,
$\spec H\subseteq\{0\}\cup[\gamma,\infty)$, and dynamics $\beta_t=\mathrm{Ad}\,e^{itH}$ with
exact finite propagation for the net $\A(\cdot)$. Let $B\subset\R$ be compact, $X=X^*$
affiliated with $\A(U)$ for every bounded open $U\supseteq B$, $\Om\in D(X)$, and
$A\in\A(O)$. If $R:=\dist(B,O)>0$ and $\gamma R\ge1$, then
\[
  \bigl|\ip{X\Om}{A\Om}-\ip{X\Om}{\Om}\ip{\Om}{A\Om}\bigr|
  \le2\,\|X\Om\|\,\|A\|\,e^{-\gamma R/2} .
\]
\end{lemma}

\begin{proof}
Set $P_0=|\Om\rangle\langle\Om|$,
$F(t)=\ip{X\Om}{(e^{itH}-P_0)A\Om}$, $G(t)=\ip{A^*\Om}{(e^{-itH}-P_0)X\Om}$: complex spectral
measures supported in $[\gamma,\infty)$, $(-\infty,-\gamma]$, variations $\le M_0=\|X\Om\|\|A\|$.
For $|t|<R$ choose a bounded open $U\supseteq B$ disjoint from $O_{|t|}$; finite propagation
puts $\beta_t(A)\in\A(O_{|t|})\subseteq\A(U)'$, and affiliation makes every element of $\A(U)'$
preserve $D(X)$ and commute with $X$ there. With $e^{-itH}\Om=\Om$,
\[
  \ip{X\Om}{e^{itH}A\Om}=\ip{\Om}{X\beta_t(A)\Om}=\ip{\Om}{\beta_t(A)X\Om}
  =\ip{A^*\Om}{e^{-itH}X\Om},
\]
and the $P_0$ terms agree since $\ip{X\Om}{\Om}\in\R$. So $F=G$ on $(-R,R)$;
Lemma~\ref{lem:spectral-main} bounds $F(0)$, which is the displayed left side.
\end{proof}

\subsection{The state derivative}

\begin{lemma}[state derivative]\label{lem:derivative-main}
Assume \textup{(G)}. In the setting of Proposition~\ref{prop:FKNpath-main}, for $0<s<1$ and
self-adjoint $A\in B(\HH)$, with $V:=V(g'-g)$:
\begin{equation}\label{eq:statederiv-main}
  \frac d{ds}\,\omega_{g_s}(A)
  =-2\,\mathrm{Re}\int_\R W_\gamma(t)\Bigl[
  \ip{V\Om_s}{\beta^{g_s}_t(A)\Om_s}-\ip{V\Om_s}{\Om_s}\,\omega_{g_s}(A)\Bigr]dt ,
\end{equation}
absolutely convergent.
\end{lemma}

\begin{proof}
By \eqref{eq:HFmain} and the parallel gauge,
$\frac d{ds}\omega_{g_s}(A)=2\mathrm{Re}\ip{\dot\Om_s}{A\Om_s}
=-2\mathrm{Re}\ip{V\Om_s}{H_s^{-1}Q_sA\Om_s}$ (all pairings defined by
Proposition~\ref{prop:standard-main}). Insert the filter
(Lemma~\ref{prop:filter-main}): $H_s^{-1}Q_sA\Om_s=\int W_\gamma(t)e^{itH_s}Q_sA\Om_s\,dt$;
since $e^{itH_s}A\Om_s=\beta^{g_s}_t(A)\Om_s$ and $P_sA\Om_s=\omega_{g_s}(A)\Om_s$, the
integrand is the bracket of \eqref{eq:statederiv-main}. Absolute convergence:
$|\ip{V\Om_s}{e^{itH_s}Q_sA\Om_s}|\le\|V\Om_s\|\|A\|$ and $W_\gamma\in L^1$.
\end{proof}

\subsection{The cell estimate and the Cauchy theorem}

Fix $\rho\in C_c^\infty((-1,1))$, $0\le\rho\le1$, $\rho>0$ on $[-\frac12,\frac12]$; put
$D(x)=\sum_k\rho(x-k)>0$ (smooth, $1$-periodic) and $\eta_j=\rho(\cdot-j)/D$: a smooth
partition of unity with $0\le\eta_j\le1$, $\supp\eta_j\subset(j-1,j+1)$, $\sum_j\eta_j=1$.
For $v=g'-g$ set $h_j=\eta_jv$, $V_j=V(h_j)$; finitely many are nonzero,
$V(v)\xi=\sum_jV_j\xi$ on common domains, and --- since $|v|\le1$, $\|h_j\|_\infty\le1$,
$\supp h_j\subset(j-1,j+1)$ --- hypothesis \textup{(M)} gives
\begin{equation}\label{eq:cellmoment-main}
  \|V_j\Om_s\|\le M\qquad(0\le s\le1).
\end{equation}

\begin{proposition}[one path]\label{prop:path-main}
Assume \textup{(G)} and \textup{(M)}. Let $O$ be a bounded interval and $g,g'\in\G$ with
$g=g'=1$ on $O_d$, $d\ge d_0:=4+2/\gamma$. Then
\begin{equation}\label{eq:pathbound-main}
  \bigl\|(\omega_g-\omega_{g'})|_{\A(O)}\bigr\|
  \ \le\ 32M\sum_{n\ge\lfloor d\rfloor}
  \Bigl(\|W_\gamma\|_1e^{-\gamma n/4}+T_\gamma(n/2)\Bigr).
\end{equation}
\end{proposition}

\begin{proof}
It suffices to bound $|\omega_g(A)-\omega_{g'}(A)|$ for self-adjoint $A\in\A(O)$, $\|A\|\le1$
($\omega_g-\omega_{g'}$ is hermitian). Since $0\le g,g'\le1$,
$g_s(x)=1\Leftrightarrow g(x)=g'(x)=1$; hence $g_s\in\G$ and $g_s=1$ on $O_d$ for all $s$, so
\textup{(G)}, \textup{(M)} and Proposition~\ref{prop:FKNpath-main} apply along the whole path.
Every nonzero $h_j$ is supported at distance $r_j:=\dist(B_j,O)\ge d$ from $O$,
$B_j:=\supp h_j$.

Fix $s\in(0,1)$, $j$, and set
$C_{j,s}(t):=\ip{V_j\Om_s}{\beta^{g_s}_t(A)\Om_s}-\ip{V_j\Om_s}{\Om_s}\omega_{g_s}(A)$.
For $|t|\le r_j/2$: Proposition~\ref{prop:standard-main}(ii) gives
$\beta^{g_s}_t(A)\in\A(O_{|t|})$ with $\dist(B_j,O_{|t|})\ge r_j/2$, and
$\gamma r_j/2\ge1$ (as $r_j\ge d\ge2/\gamma$); Lemma~\ref{prop:unbdd-main} with $X=V_j$
(affiliation and $\Om_s\in D(V_j)$ by \S\ref{sec:setting} and
Proposition~\ref{prop:standard-main}(i)) and \eqref{eq:cellmoment-main} give
\begin{equation}\label{eq:near-main}
  |C_{j,s}(t)|\le2Me^{-\gamma r_j/4}.
\end{equation}
For all $t$, Cauchy--Schwarz gives $|C_{j,s}(t)|\le2M$. Substituting $V=\sum_jV_j$ into
\eqref{eq:statederiv-main} and splitting each $t$-integral at $|t|=r_j/2$,
\[
  \Bigl|\frac d{ds}\omega_{g_s}(A)\Bigr|
  \le4M\sum_j\Bigl(\|W_\gamma\|_1e^{-\gamma r_j/4}+T_\gamma(r_j/2)\Bigr).
\]
\emph{Shell count.} Write $O=(a,b)$. Since $d\ge4$ and $B_j\subset[j-1,j+1]$, a nonzero $B_j$
lies on one side of $O_d$; if on the right with $r_j\in[n,n+1)$, compactness gives $x_j\in B_j$
with $x_j-b=r_j$, and $|j-x_j|\le1$ forces $j\in[b+n-1,b+n+2]$ --- at most four integers; same
on the left: at most eight $j$ per shell. Both $r\mapsto e^{-\gamma r/4}$ and
$r\mapsto T_\gamma(r/2)$ are decreasing and $r_j\ge d$, so the sum over $j$ is at most
$8\sum_{n\ge\lfloor d\rfloor}(\cdots)$, giving the bound $32M\sum_{n\ge\lfloor d\rfloor}(\cdots)$
for the derivative. Integrate over $[\varepsilon,1-\varepsilon]$ and use the endpoint norm
continuity of $P_s$ (Proposition~\ref{prop:FKNpath-main}), then $\varepsilon\downarrow0$ and the
supremum over the unit ball.
\end{proof}

Define
\begin{equation}\label{eq:Psi-main}
 \Psi_\gamma(d):=
 \begin{cases}
  \displaystyle 32M\sum_{n\ge\lfloor d\rfloor}
  \Bigl(\|W_\gamma\|_1e^{-\gamma n/4}+T_\gamma(n/2)\Bigr),&d\ge d_0,\\[6pt]
  \max\{2,\Psi_\gamma(d_0)\},&0\le d<d_0.
 \end{cases}
\end{equation}

\begin{lemma}[decay of the rate function]\label{lem:Psitail-main}
For every $N\ge1$ there is $C_{N,\gamma,M}$ with $\Psi_\gamma(d)\le C_{N,\gamma,M}(1+d)^{-N}$.
\end{lemma}

\begin{proof}
The exponential series has an exponentially decaying tail; for the other,
$T_\gamma(r)\le r^{-q}\int|t|^q|W_\gamma|$ (Lemma~\ref{prop:filter-main}), take $q=N+2$ and sum
$n^{-q}$.
\end{proof}

\begin{theorem}[Cauchy/LPPL]\label{thm:lppl-main}
Assume \textup{(G)} and \textup{(M)} --- supplied here by Theorems~\ref{thm:G-main} and
\ref{thm:M2-main} with $\gamma=m_1$. Then there is a state $\omega_\infty$ on $\A$ with
$\lim_{g\in\G}\|(\omega_g-\omega_\infty)|_{\A(O)}\|=0$ for every bounded interval $O$;
for $g\equiv1$ on $O_d$,
\[
  \|(\omega_g-\omega_\infty)|_{\A(O)}\|\ \le\ \Psi_{m_1}(d);
\]
and $\omega_\infty|_{\A(O)}$ is normal for every $O$.
\end{theorem}

\begin{proof}
Fix $O$, $\varepsilon>0$; by Lemma~\ref{lem:Psitail-main} pick $d\ge d_0$ with
$\Psi_{m_1}(d)<\varepsilon$, and $g_0\in\G$ with $g_0\equiv1$ on $O_d$. If
$g,g'\succcurlyeq g_0$ then both equal $1$ on $O_d$ (in particular on a common interval), so
Proposition~\ref{prop:path-main} applies directly to the pair and gives
$\|(\omega_g-\omega_{g'})|_{\A(O)}\|\le\Psi_{m_1}(d)<\varepsilon$. Hence
$(\omega_g|_{\A(O)})_{g\in\G}$ is Cauchy in $\A(O)^*$; the limits are compatible under isotony
and extend by norm continuity to a state $\omega_\infty$ on $\A$. The rate follows by fixing
$g\equiv1$ on $O_d$ and letting $g'$ run through the cofinal tail of cutoffs equal to $1$ on
$O_d$; for $d<d_0$ the bound is trivial. Normality: normal functionals form a norm-closed
subspace of $\A(O)^*$ --- if $0\le A_i\uparrow A$ in $\A(O)$ and
$\|(\omega_g-\omega_\infty)|_{\A(O)}\|<\delta$, then
$0\le\omega_\infty(A)-\omega_\infty(A_i)\le2\delta\|A\|+\omega_g(A-A_i)$; let $i\uparrow$ (normality
of $\omega_g$), then $\delta\downarrow0$.
\end{proof}

\section{The limit is a ground state; invariances}\label{sec:ground}

We import from companion M0 (proofs there are complete and short):

\begin{proposition}[exact spectral passage; M0]\label{prop:M0-main}
\emph{(i)} If $(g_n)\subseteq\G$ satisfies $d(g_n,O)\to\infty$ for every bounded $O$, and
$\omega_{g_n}\to\omega$ in norm on every $\A(O)$, then $\omega$ is a ground state of
$(\A,\alpha)$ in the spectral sense: $\omega\circ\alpha_t=\omega$, the function
$t\mapsto\omega(A\alpha_t(B))$ is continuous and bounded for all $A,B\in\A$, and
$\int f(t)\,\omega(A\alpha_t(B))\,dt=0$ for every $f\in L^1(\R)$ whose transform
$\check f(E)=\int f(t)e^{itE}dt$ vanishes on $[0,\infty)$.
\emph{(ii)} $\omega_{\tau_ag}=\omega_g\circ\tau_{-a}$ for all $a\in\R$, $g\in\G$.
\emph{(iii)} For even $\mathscr P$, $\omega_g\circ\theta=\omega_g$ for every $g$, where
$\theta=\mathrm{Ad}\,\Gamma(-1)$.
\end{proposition}

\begin{proof}
\emph{(ii)} $U_aK(g)U_a^*=K(\tau_ag)$ (the free part commutes with translations; the
interaction transforms covariantly), so by simplicity (Source Fact~\ref{fact:sa})
$\Om_{\tau_ag}=U_a\Om_g$ up to a phase, which cancels in the state.

\emph{(iii)} $\Theta:=\Gamma(-1)$ is unitary, self-adjoint, commutes with $H_0$, and for even
$\mathscr P$ satisfies $\Theta V(g)\Theta=V(g)$; hence $\Theta\Om_g=c\Om_g$ with $c\in\{\pm1\}$
by simplicity and $\Theta^2=1$. In $Q$-space $(\Theta\psi)(q)=\psi(-q)$, so $\Theta\Om_g>0$
a.e.\ together with $\Om_g>0$ forces $c=+1$.

\emph{(i)} \emph{Exact identity.} Fix a bounded interval $O$, $A,B\in\A(O)$,
$M_0:=\|A\|\|B\|$, and $g\in\G$; let
$\rho_g(\cdot):=\ip{A^*\Om_g}{P_{H(g)}(\cdot)B\Om_g}$, a complex measure with
$\supp\rho_g\subseteq\spec H(g)\subseteq[0,\infty)$ and $|\rho_g|(\R)\le M_0$, and
$G_g(t):=\int e^{itE}d\rho_g(E)$. For $|t|<d(g,O)$, the cutoff equals one on $O_{|t|}$, so by
\eqref{eq:P} and $e^{-itH(g)}\Om_g=\Om_g$,
\[
  \omega_g\bigl(A\alpha_t(B)\bigr)
  =\ip{\Om_g}{Ae^{itH(g)}Be^{-itH(g)}\Om_g}
  =\ip{A^*\Om_g}{e^{itH(g)}B\Om_g}=G_g(t) .
\]
Moreover $\int_\R f\,G_g\,dt=\int_{[0,\infty)}\check f\,d\rho_g=0$ for the stated class of $f$
(Fubini; $\check f=0$ on the support).

\emph{Invariance.} For $A\in\A(O)$ and fixed $t$, take $n$ so large that
$d(g_n,O)>|t|$. Then \eqref{eq:P} and $e^{-itH(g_n)}\Om_{g_n}=\Om_{g_n}$ give directly
\[
 \omega_{g_n}(\alpha_t(A))
 =\ip{\Om_{g_n}}{e^{itH(g_n)}Ae^{-itH(g_n)}\Om_{g_n}}
 =\omega_{g_n}(A).
\]
Both $\alpha_t(A)$ and $A$ belong to $\A(O_{|t|})$, so norm convergence on
$\A(O_{|t|})$ passes to the limit:
$\omega(\alpha_t(A))=\omega(A)$, and by density on $\A$.

\emph{Continuity and vanishing.} Fix $A,B\in\A(O)$ and $T>0$. For $n$ with $d(g_n,O)>T$ and
$|t|\le T$: $A\alpha_t(B)\in\A(O_T)$ with norm $\le M_0$, so
\[
  \sup_{|t|\le T}\bigl|\omega(A\alpha_t(B))-\omega_{g_n}(A\alpha_t(B))\bigr|
  \le M_0\,\bigl\|(\omega-\omega_{g_n})|_{\A(O_T)}\bigr\|\longrightarrow0 :
\]
$t\mapsto\omega(A\alpha_t(B))$ is a locally uniform limit of the continuous functions $G_{g_n}$,
hence continuous, bounded by $M_0$. For $f$ as stated and
$\varepsilon(T):=\int_{|t|>T}|f|$:
\[
  \int_\R f\,\omega(A\alpha_\cdot(B))
  =\underbrace{\int_{|t|\le T}f\,[\omega-\omega_{g_n}](A\alpha_t(B))}_{\le\|f\|_1M_0
  \|(\omega-\omega_{g_n})|_{\A(O_T)}\|}
  +\underbrace{\int_{|t|\le T}f\,G_{g_n}}_{=\,-\int_{|t|>T}f\,G_{g_n}}
  +\int_{|t|>T}f\,\omega(A\alpha_t(B)) ,
\]
so its modulus is $\le\|f\|_1M_0\|(\omega-\omega_{g_n})|_{\A(O_T)}\|+2M_0\varepsilon(T)$;
let $n\to\infty$, then $T\to\infty$. Joint norm continuity in $(A,B)$ extends the conclusion
from local $A,B$ to $\A$.
\end{proof}

\begin{theorem}[Main Theorem]\label{thm:MAIN}
Let $0\le\lambda/m_0^2<\varepsilon_{\mathscr P}$. Then there is a state $\omega_\infty$ on $\A$
such that:
\begin{enumerate}
\item[\rm(1)] \emph{(full-net convergence with rate)} for every bounded open interval $O$,
      $\lim_{g\in\G}\|(\omega_g-\omega_\infty)|_{\A(O)}\|=0$, and
      $\|(\omega_g-\omega_\infty)|_{\A(O)}\|\le\Psi_{m_1}(d)$ whenever $g\equiv1$ on $O_d$,
      with $\Psi_{m_1}$ decaying faster than every inverse power;
\item[\rm(2)] \emph{(local normality)} $\omega_\infty|_{\A(O)}$ is normal for every $O$;
\item[\rm(3)] \emph{(ground state)} $\omega_\infty$ is a ground state of $(\A,\alpha)$ in the
      spectral sense;
\item[\rm(4)] \emph{(invariances)} $\omega_\infty\circ\tau_a=\omega_\infty$ for all $a\in\R$;
      and for even $\mathscr P$, $\omega_\infty\circ\theta=\omega_\infty$;
\item[\rm(5)] \emph{(spectra and local moment)}
      \[
       \spec H(g)\subseteq\{0\}\cup[m_1,\infty)\quad(g\in\G),
       \qquad \dim\ker H(g)=1,
      \]
      and $\sup_g\|V(h)\Om_g\|\le M$ for unit-bounded $h$ supported in length-$2$ intervals.
\end{enumerate}
The only nonstandard weak-coupling estimates imported are the printed GJS
Theorems 1.1.7--1.1.8 (Source Facts \ref{sf:CE2}--\ref{sf:CE1}, with the cutoff uniformity of
Remark~\ref{rem:audit}). In addition, the proof uses the separately identified standard
cutoff-construction inputs: self-adjointness and the simple positive cutoff ground state
(Source Fact~\ref{fact:sa}); FKN transfer, finite-volume stability, and free
hypercontractivity (Lemma~\ref{lem:FKN-main} and Proposition~\ref{prop:FKNpath-main});
model-specific domain smoothing and the common core
(Propositions~\ref{prop:standard-main} and \ref{prop:FKNpath-main}); finite propagation and
cutoff patching; and translation/parity covariance (Source Fact~\ref{fact:PT} and
Proposition~\ref{prop:M0-main}). Proposition~\ref{prop:scaling} verifies that the dimensionless
weak-coupling hypothesis implies the source smallness condition.
\end{theorem}

\begin{proof}
(5) is Theorems~\ref{thm:G-main} and \ref{thm:M2-main}. These supply (G), (M) with
$\gamma=m_1$, so Theorem~\ref{thm:lppl-main} gives (1)--(2). For (3): choose
$\widetilde g_n\in\G$ with $\widetilde g_n\equiv1$ on $(-n,n)$; then $d(\widetilde g_n,O)\to\infty$
for every $O$ and, by (1), $\omega_{\widetilde g_n}\to\omega_\infty$ in norm on every $\A(O)$;
Proposition~\ref{prop:M0-main}(i) applies. For (4): translation is an order isomorphism of
$\G$ with cofinal range, so by \eqref{eq:T} and (1),
$\omega_\infty(A)=\lim_g\omega_{\tau_ag}(A)=\lim_g\omega_g(\tau_{-a}A)
=\omega_\infty(\tau_{-a}A)$ for local $A$, and by density on $\A$. The $\Z_2$ statement is the
local norm limit of Proposition~\ref{prop:M0-main}(iii).
\end{proof}

\begin{corollary}[the Glimm--Jaffe vacuum is canonical at weak coupling]\label{cor:canonical}
Every $w^*$ limit point of the (averaged or unaveraged) finite-volume ground states constructed
in \cite[\S2]{GJ3} coincides on each $\A(O)$ with $\omega_\infty$; in particular the
construction of \cite{GJ3} is subsequence- and averaging-independent in the weak-coupling
region.
\end{corollary}

\begin{proof}
For the unaveraged net this is Theorem~\ref{thm:MAIN}(1). For the averaged states of
\cite[(2.3)--(2.5)]{GJ3}: with $g_n=g(\cdot/n)$, $g\equiv1$ on $[-3,3]$, and $h\ge0$,
$\supp h\subseteq[-1,1]$, $\int h=1$,
\begin{align*}
  \omega_n(A)
  &=\frac1n\int\ip{\Om_{g_n}}{\sigma_\alpha(A)\Om_{g_n}}\,
      h(\alpha/n)\,d\alpha\\
  &=\int_{-1}^{1}h(u)\,\omega_{g_n}\bigl(\tau_{\epsilon nu}A\bigr)\,du\\
  &=\int_{-1}^{1}h(u)\,\omega_{\tau_{-\epsilon nu}g_n}(A)\,du ,
\end{align*}
by \eqref{eq:T}, where $\epsilon\in\{\pm1\}$ fixes the direction convention of $\sigma_\alpha$
(the estimate below is invariant under $\epsilon\to-\epsilon$). Since $g_n\equiv1$ on
$[-3n,3n]$, the shifted cutoff $\tau_{-\epsilon nu}g_n$ equals $1$ on an interval containing
$[-2n,2n]$ for every $|u|\le1$; hence for $O\subseteq[-R,R]$,
$d(\tau_{-\epsilon nu}g_n,O)\ge2n-R$ \emph{uniformly in $u$}, and Theorem~\ref{thm:MAIN}(1)
gives
\[
  \|(\omega_n-\omega_\infty)|_{\A(O)}\|
  \le\int_{-1}^1h(u)\,\bigl\|(\omega_{\tau_{-\epsilon nu}g_n}-\omega_\infty)|_{\A(O)}\bigr\|\,du
  \le\Psi_{m_1}(2n-R)\longrightarrow0 .
\]
No translation invariance of $\omega_\infty$ is needed for this bound.
\end{proof}

\begin{corollary}[conditional exclusion of a two-phase structure]\label{cor:nobreak}
Assume that $\mathscr P$ is even and $0\le\lambda/m_0^2<\varepsilon_{\mathscr P}$. The following
four assertions cannot hold simultaneously:
\begin{enumerate}
\item[\rm(B1)] there are translation-invariant ground states $\omega_+$ and $\omega_-$ of
      $(\A,\alpha)$ such that, for all bounded local $X,Y$,
      \[
       \omega_\pm(X\tau_a(Y))-\omega_\pm(X)\omega_\pm(\tau_a(Y))
       \longrightarrow0\qquad(|a|\to\infty);
      \]
\item[\rm(B2)] $\omega_+\circ\theta=\omega_-$ and $\omega_-\circ\theta=\omega_+$;
\item[\rm(B3)] there are a bounded local observable $B=B^*\in\A(O)$ and $m\ne0$ such that
      $\omega_\pm(B)=\pm m$;
\item[\rm(B4)] every $\theta$-invariant ground state of $(\A,\alpha)$ equals
      $\frac12(\omega_++\omega_-)$.
\end{enumerate}
\end{corollary}

\begin{proof}
Assume (B1)--(B4). Choose $\widetilde g_n\in\G$ with
$\widetilde g_n\equiv1$ on $(-n,n)$ and write $\omega_n=\omega_{\widetilde g_n}$. Fix
$a\in\R$ such that
\[
 R_a:=\dist(O,O+a)>0,\qquad m_1R_a\ge1.
\]
For the cutoff Hamiltonian $H(\widetilde g_n)$ apply Lemma~\ref{prop:unbdd-main} with
unbounded-argument slot $X=B$, bounded-argument slot $A=\tau_a(B)$, localization set
$\overline O$, and $\gamma=m_1$. The bounded operator $B$ has domain $\HH$; because
$B\in\A(O)\subseteq\A(U)$ for every bounded open $U\supseteq\overline O$, it is affiliated
with every such $\A(U)$. Thus all hypotheses of that lemma hold. Since
$\|B\Om_{\widetilde g_n}\|\le\|B\|$ and $\|\tau_a(B)\|=\|B\|$, it gives, uniformly in $n$,
\begin{equation}\label{eq:cutoff-bounded-cluster}
 \left|\omega_n(B\tau_a(B))-\omega_n(B)\omega_n(\tau_a(B))\right|
 \le2\|B\|^2e^{-m_1R_a/2}.
\end{equation}
For fixed $a$, all three observables in \eqref{eq:cutoff-bounded-cluster} belong to the local
algebra of one bounded interval containing $O\cup(O+a)$. The local-norm convergence of
Theorem~\ref{thm:MAIN}(1) therefore permits passage to $n\to\infty$. Translation invariance
of $\omega_\infty$ then yields
\[
 \left|\omega_\infty(B\tau_a(B))-\omega_\infty(B)^2\right|
 \le2\|B\|^2e^{-m_1R_a/2}\xrightarrow[|a|\to\infty]{}0.
\]
By Theorem~\ref{thm:MAIN}(3)--(4), $\omega_\infty$ is a $\theta$-invariant ground state.
Assumption (B4) consequently gives
\[
 \omega_\infty=\tfrac12(\omega_++\omega_-),\qquad
 \omega_\infty(B)=\tfrac12(m-m)=0.
\]
By (B1), translation invariance and clustering of each $\omega_\pm$ imply
\[
 \omega_\pm(B\tau_a(B))
 -\omega_\pm(B)\omega_\pm(\tau_a(B))\longrightarrow0,\qquad
 \omega_\pm(\tau_a(B))=\omega_\pm(B)=\pm m.
\]
Hence
\[
 \omega_\infty(B\tau_a(B))
 =\tfrac12\!\left(\omega_+(B\tau_a(B))+\omega_-(B\tau_a(B))\right)
 \longrightarrow\tfrac12(m^2+m^2)=m^2\ne0,
\]
contradicting the preceding cluster bound, whose left side tends to zero because
$\omega_\infty(B)=0$.
\end{proof}

\section{Discussion}\label{sec:discussion}

\emph{Relation to GRS Problems 1--2.} Problem 2 of \cite[p.~126]{GRS75} asks for local
convergence of Euclidean Gibbs states; Theorem~\ref{thm:MAIN}(1) is a Hamiltonian counterpart,
in norm on local von Neumann algebras, for the vacuum family. Problem 1 asks for uniform local
$L^p$ bounds on the Euclidean vacuum density. The note added on \cite[p.~128]{GRS75} records
Newman's solution of the small-coupling $p=1$ versions, so the printed problems must not be read
as an unqualified current open-problem claim. The present proof neither establishes nor needs
Problem 1: the identification step uses the exact spectral identity of
Proposition~\ref{prop:M0-main}, and the convergence step uses the Cauchy estimate.

\emph{Sharpness.} In the quadratic model the exact rate is $\mathrm{poly}(d)\,e^{-2m_1d}$
(companion M2: round trip through the region where the Hamiltonians differ, at the local Agmon
mass). $\Psi_{m_1}$ is super-polynomial but sub-exponential. The present filter argument does
not obtain the sharp exponential rate; a proof of such a rate would require an additional
resolvent or Combes--Thomas input not supplied here.

\emph{Uniqueness.} $\omega_\infty$ is canonical for the cutoff family, and translation
invariant without averaging. The theorem does not identify every locally normal
translation-invariant ground state of $(\A,\alpha)$ with $\omega_\infty$. In particular, the
Euclidean results \cite{AHZ,BDW} do not by themselves provide the reverse reconstruction and
state-space identification needed for that Hamiltonian conclusion; companion M3 records that
missing implication as (R3).

\emph{Beyond weak coupling.} Corollary~\ref{cor:nobreak} is conditional on (B1)--(B4).
Under exactly those hypotheses, a proof based on a cutoff-uniform gap cannot extend to that
two-phase regime. The phase-transition results \cite{GJS75,GJS76}, as cited here, are not used
as a substitute for (B4), and this manuscript makes no unconditional assertion that its
Hamiltonian state space has the classification (B1)--(B4).

\emph{Part II.} Uniformity in a wavelet resolution --- the OAR programme --- is a
separate problem.  For the MMST free nearest-neighbour model and its Daubechies
observable embedding \cite{MMST}, the companion Lamport proof constructs local Weyl
unitaries for which
\[
 \liminf_{N\to\infty}
 \|\omega_N-\omega_\infty\mathbin\circ\beta_N\|_{\mathfrak W_N(O)}
 \ge e^{-\sqrt{2+\sqrt2}/4}-e^{-3\pi/16}>0.
\]
Thus no diagonal local-state-norm error $C\varepsilon_N^\theta$, $\theta>0$, is possible.
The implemented replacement fixes a torus $\mathbb T_L$, the MMST nearest-neighbour
Hamiltonian, the exact $D4$ observable maps, and the finite-torus Wick convention.
The finite-resolution interacting Schr\"odinger form is closed and lower bounded; its
operator has compact resolvent and a simple ground state.

For these interacting ground states, weak-star compactness gives a cofinal subnet
$N(i)$ and a state $\Omega_{L,\lambda}^{\mathrm{cl}}$ on the inductive-limit algebra such
that, for every fixed $M$ and $A\in\mathfrak W_{M,L}$,
\[
 \lim_i\omega_{N(i),L,\lambda,1}
   \bigl(\alpha_M^{N(i)}(A)\bigr)
 =\Omega_{L,\lambda}^{\mathrm{cl}}\bigl(\alpha_M^\infty(A)\bigr).
\]
This proves existence of a projective interacting cluster state, not convergence of the
full sequence or its identification with the continuum vacuum.  In the free case the
full sequence satisfies the explicit estimate
\[
\begin{split}
 &\left|\omega_{M+r,L}^{(0)}\!\circ\alpha_M^{M+r}(W_M(\xi))
 -\omega_L^{(0)}\!\circ\beta_M(W_M(\xi))\right|\\
 &\hspace{35mm}\le\frac{C_{M,L,m,\xi}}4\,2^{-\theta r},
 \qquad
 \theta=\frac{4(1-\log_2\sqrt3)}{7-2\log_2\sqrt3}>0.
\end{split}
\]

For the static interacting theorem, the exact remaining assertion is
\[
 \lim_{N\to\infty}
 \omega_{N,L,\lambda,1}\bigl(\alpha_M^N(W_M(\xi))\bigr)
 =
 \omega_{P(\phi)_2,L}\bigl(\beta_M(W_M(\xi))\bigr)
 \quad\text{for every fixed }M,\xi.                              \tag{PII}\label{eq:PII}
\]
The companion Weyl criterion proves that \eqref{eq:PII}, by itself, extends to every
$A\in\mathfrak W_{M,L}$, gives projective consistency, and identifies the
inductive-limit state.  A uniform gap, local second moment, and interacting
Lieb--Robinson bound are therefore not prerequisites for this static fixed-torus
criterion.  They remain necessary for stronger quantitative-locality or dynamical
claims.  In particular, Theorem~4.1 of \cite{NRSS} assumes $V'\in L^1(\mathbb R)$
and cannot be applied to the Wick quartic, since
$|V_c'(u)|\ge2\lambda|u|^3$ for $|u|\ge\sqrt{6c}$.
No full-sequence identified interacting resolution-limit theorem is asserted here.
The proofs and exact ledgers are in \texttt{strategy.md} and the five Lamport documents
in \texttt{part\_ii/}.

\appendix
\section{Interface checklist}\label{app:checklist}

\begin{itemize}
\item (G), (M) of \texttt{lppl.tex}: supplied with $\gamma=m_1$, $M$ as in
      Theorem~\ref{thm:M2-main}. Hypothesis classes match ($\G$; $\|h\|_\infty\le1$, support in
      length-$2$ intervals). $\checkmark$
\item \eqref{eq:P}, \eqref{eq:T}: supplied by Source Fact~\ref{fact:PT}. $\checkmark$
\item Simplicity of $E(g)$ (used in Lemma~\ref{lem:vac-main} and
      Prop.~\ref{prop:M0-main}(iii)): Source Fact~\ref{fact:sa}. No isolation input anywhere.
      $\checkmark$
\item $\Om_g\in D(V(h))$ (used in Prop.~\ref{prop:unbdd-main}):
      Prop.~\ref{prop:standard-main}(i), and independently re-proved in
      Theorem~\ref{thm:M2-main} Step 3. $\checkmark$
\item Wick/covariance conventions between \cite{GJS} and \S\ref{sec:setting}:
      Remark~\ref{rem:conv}. $\checkmark$
\item Gate R2: \emph{closed by page-image inspection} (\texttt{m4/r2-check.tex}) ---
      $K_1$ to the first power, local weight $K_2^{2N(\Delta)}$; all displayed constants use
      the exact reading. $\checkmark$
\end{itemize}

\begin{thebibliography}{99}
\bibitem{GJ1} J.~Glimm, A.~Jaffe, Phys.\ Rev.\ \textbf{176} (1968) 1945--1951.
\bibitem{GJ2} J.~Glimm, A.~Jaffe, Ann.\ of Math.\ \textbf{91} (1970) 362--401.
\bibitem{GJ3} J.~Glimm, A.~Jaffe, Acta Math.\ \textbf{125} (1970) 203--267.
\bibitem{GJ4} J.~Glimm, A.~Jaffe, J.\ Math.\ Phys.\ \textbf{13} (1972) 1568--1584.
\bibitem{GJ71} J.~Glimm, A.~Jaffe, Comm.\ Math.\ Phys.\ \textbf{22} (1971) 253--258.
\bibitem{GJExp} J.~Glimm, A.~Jaffe, \emph{Quantum Field Theory and Statistical Mechanics:
Expositions}, Birkh\"auser, 1985.
\bibitem{GJS} J.~Glimm, A.~Jaffe, T.~Spencer, Ann.\ of Math.\ \textbf{100} (1974) 585--632.
\bibitem{GJS75} J.~Glimm, A.~Jaffe, T.~Spencer, Comm.\ Math.\ Phys.\ \textbf{45} (1975) 203--216.
\bibitem{GJS76} J.~Glimm, A.~Jaffe, T.~Spencer, Ann.\ Physics \textbf{101} (1976) 610--669.
\bibitem{GRS75} F.~Guerra, L.~Rosen, B.~Simon, Ann.\ of Math.\ \textbf{101} (1975) 111--259.
\bibitem{Rosen70} L.~Rosen, Comm.\ Math.\ Phys.\ \textbf{16} (1970) 157--183.
\bibitem{SimonBook} B.~Simon, \emph{The $P(\Phi)_2$ Euclidean (Quantum) Field Theory},
Princeton UP, 1974.
\bibitem{AHZ} S.~Albeverio, R.~H\o egh-Krohn, B.~Zegarlinski, Comm.\ Math.\ Phys.\ \textbf{123}
(1989) 377--424.
\bibitem{BDW} R.~Bauerschmidt, B.~Dagallier, H.~Weber, arXiv:2504.08606.
\bibitem{RZT} M.~Rodriguez Zarate, T.~Thiemann, arXiv:2505.13030.
\bibitem{MMST} V.~Morinelli, G.~Morsella, A.~Stottmeister, Y.~Tanimoto,
Commun. Math. Phys. \textbf{387} (2021) 1299--1367.
\bibitem{NRSS} B.~Nachtergaele, H.~Raz, B.~Schlein, R.~Sims,
Commun. Math. Phys. \textbf{286} (2009) 1073--1098.
\end{thebibliography}

\end{document}
